\documentclass[11pt]{article}
\usepackage[margin=1in]{geometry}
\usepackage{amsmath,amssymb,bm}
\usepackage{booktabs}
\usepackage{hyperref}

\title{0D--2P Relativistic Electron Fokker--Planck Model\\Physical and Numerical Reference}
\author{}
\date{September 2026}

\newcommand{\p}{p}
\newcommand{\x}{\xi}
\newcommand{\g}{\gamma}
\newcommand{\N}{N}
\newcommand{\f}{f}
\newcommand{\Ebar}{\overline{E}}

\begin{document}
\maketitle

\noindent
This document specifies the spatially homogeneous 0D--2P forward relativistic
electron Fokker--Planck model and its finite-volume realization.  Standalone mode
includes electric-field acceleration, relativistic small-angle collisions,
synchrotron radiation reaction, and an optional Chiu--Harvey large-angle
secondary-electron source.  Standalone mode uses a prescribed constant plasma
state.  A bulk-plasma driver is implemented and exchanges stage plasma states and
kinetic current moments through a monolithic analytic-Newton solve while preserving
the conservative finite-volume and direct-solver contracts defined here.

The numerical reference covers the dimensional and normalized equations, fixed
coefficient construction, conservative cell-content discretization, boundary
conditions, implicit time integration, sparse CSR assembly, and PETSc direct solves,
bulk-plasma coupling, and numerical qualification.  Its scope includes forward
evolution, coupled evolution, and diagnostics required to qualify both paths.

Revision v16.0 is the living physical and numerical reference for current forward
and coupled implementations.  The equations and numerical conventions below are
authoritative for the maintained solver, physics, data, and TOML files.

%===============================================================================
\section{Physical model}

The present model is a spatially homogeneous, gyrotropic electron kinetic system in
momentum magnitude and pitch angle.  Standalone mode prescribes all plasma quantities
and holds them constant in time.  Coupled mode evolves a host-side bulk state and the
kinetic state together at every TR--BDF2 stage; spatial transport remains outside the
current model.  The physical system is defined below in dimensional form.  No reference
normalization used by the numerical solver is introduced until the final subsection of
this section.

\subsection{Dimensional kinetic equation}

The physical model is defined first in dimensional variables.  Let
$F_e(P,\xi,t)$ be the gyrophase-averaged electron distribution function, where $P$ is
the physical momentum magnitude and
\begin{equation}
 \xi \equiv \frac{P_\parallel}{P}\in[-1,1].
\end{equation}
The forward kinetic equation is
\begin{equation}
 \frac{\partial F_e}{\partial t}
 = \mathcal C_{\rm sa}[F_e]
 + \mathcal E[F_e]
 + \mathcal R_{\rm syn}[F_e]
 + \mathcal S_{\rm LA}[F_e].
 \label{eq:dimensional_kinetic}
\end{equation}
Here $\mathcal C_{\rm sa}$ is the relativistic small-angle test-particle
Fokker--Planck collision operator, $\mathcal E$ is acceleration by the parallel
electric field, $\mathcal R_{\rm syn}$ is synchrotron radiation reaction, and
$\mathcal S_{\rm LA}$ is either zero or the Chiu--Harvey large-angle
secondary-electron gain operator defined below.  The same forward operator structure is standard in relativistic runaway-electron
kinetics; see, for example, Ref.~\cite{BraamsKarney1989}.

The dimensional distribution is defined such that its momentum-space integral gives the
represented electron density,
\begin{equation}
 n_F(t)
 = 2\pi\int_0^\infty P^2\,dP\int_{-1}^{1}F_e(P,\xi,t)\,d\xi,
 \label{eq:dimensional_density}
\end{equation}
so that $n_F$ is the represented electron density.  The relativistic Lorentz factor
and speed are
\begin{equation}
 \gamma(P)=\sqrt{1+\frac{P^2}{m_e^2c^2}},
 \qquad
 v(P)=\frac{P}{\gamma m_e},
 \qquad
 \beta(P)=\frac{v}{c}.
 \label{eq:dimensional_kinematics}
\end{equation}

\subsection{Parallel electric-field operator}

For an electron of charge $-e$ in a prescribed uniform parallel electric field
$E_\parallel$, the momentum-space characteristics are
\begin{equation}
 \dot P_E=-eE_\parallel\xi,
 \qquad
 \dot\xi_E=-\frac{eE_\parallel}{P}(1-\xi^2).
 \label{eq:electric_characteristics}
\end{equation}
The corresponding conservative kinetic operator is
\begin{equation}
 \mathcal E[F_e]
 =-\frac{1}{P^2}\frac{\partial}{\partial P}
   \left(P^2\dot P_E F_e\right)
  -\frac{\partial}{\partial\xi}
   \left(\dot\xi_E F_e\right).
 \label{eq:electric_operator_dimensional}
\end{equation}
This is the standard homogeneous relativistic electric-field operator used in
runaway-electron Fokker--Planck models; see
Refs.~\cite{ConnorHastie1975}.

\subsection{Small-angle test-particle collision-model hierarchy}

The standalone kinetic model admits two test-particle approximations for the
small-angle collision kernel.  The current CPU prototype uses the fully-relativistic
and asymptotically matched model historically used in CODE.  The finite-temperature
relativistic model defined first below remains the preferred reference formulation for
future implementation.  These define two alternative
\emph{thermal-log complete-screening base kernels}.  The physical coefficient construction
then applies two common layers that are independent of which base kernel is selected:
(i) energy-dependent Coulomb-logarithm matching of the electron--electron and
electron--ion sectors, and (ii) partial-screening corrections from bound electrons.
Neither layer modifies radial energy diffusion.  Thus changing the base kernel does not
change the Coulomb-logarithm or screening prescription.  Section~2 specifies the same
layered construction in the common solver normalization, without changing the grid, time
integrator, or sparse-solver architecture.

\paragraph{Finite-temperature relativistic test-particle model (reference formulation).}
Small-angle collisions are represented by a linearized relativistic test-particle
Fokker--Planck operator with radial drag, radial energy diffusion, and pitch-angle
deflection,
\begin{equation}
 \mathcal C_{\rm sa}[F_e]
 =\frac{1}{P^2}\frac{\partial}{\partial P}
  \left[
   P^2\left(
      C_F^{(P)} F_e
      + C_A^{(P)}\frac{\partial F_e}{\partial P}
   \right)
  \right]
 +\frac{\nu_D}{2}\frac{\partial}{\partial\xi}
  \left[(1-\xi^2)\frac{\partial F_e}{\partial\xi}\right].
 \label{eq:collision_operator_dimensional}
\end{equation}
The dimensional radial coefficients are written in terms of the slowing-down and
parallel diffusion frequencies as
\begin{equation}
 C_F^{(P)}=P\nu_s,
 \qquad
 C_A^{(P)}=\frac{P^2}{2}\nu_\parallel.
 \label{eq:dimensional_collision_coeffs}
\end{equation}
Equation~\eqref{eq:collision_operator_dimensional} is the relativistic test-particle
operator of Braams--Karney/Pike--Rose form.  Its thermal-log complete-screening base
coefficients are written below; the common energy-dependent Coulomb-logarithm and
partial-screening layers are then specified separately following
Refs.~\cite{BraamsKarney1989,PikeRose2014,Hesslow2018Collision}.
The field-particle part of the small-angle collision operator is not retained in the
present model.

The published collision coefficients are most compactly written using the auxiliary
relativistic kinematic variables
\begin{equation}
 u\equiv\frac{P}{m_ec},
 \qquad
 \gamma=\sqrt{1+u^2},
 \qquad
 \Theta\equiv\frac{T_e}{m_ec^2}.
 \label{eq:collision_auxiliary_variables}
\end{equation}
These dimensionless quantities are only arguments used to write the dimensional collision
frequencies compactly; they do not constitute the reference normalization of the kinetic
PDE used by the numerical solver.  Define the physical relativistic collision time
\begin{equation}
 \tau_c
 =\frac{4\pi\epsilon_0^2m_e^2c^3}
 {n_e e^4\ln\Lambda_0},
 \label{eq:physical_collision_time}
\end{equation}
where $n_e$ and $T_e$ are the prescribed free-electron density and temperature and
$\ln\Lambda_0$ is the thermal Coulomb logarithm.  Following the finite-temperature
construction in Ref.~\cite{Hesslow2018Collision}, first define the finite-temperature
\emph{thermal-log} complete-screening base by setting both collisional logarithms equal
to $\ln\Lambda_0$.  It is useful to keep the electron--electron and electron--ion
parts of the deflection frequency separate:
\begin{align}
 \nu_s^{\rm FT,0}
 &=\frac{1}{\tau_c}\frac{\gamma^2}{u^3}
   \frac{u^2}{\gamma^2}\Psi_s(u,\Theta),
 \label{eq:physical_nus}\\
 \nu_\parallel^{\rm FT,0}
 &=\frac{1}{\tau_c}\frac{2\gamma\Theta}{u^3}\Psi_s(u,\Theta),
 \label{eq:physical_nupar}\\
 \nu_{D,ee}^{\rm FT,0}
 &=\frac{1}{\tau_c}\frac{\gamma}{u^3}
   \frac{2u}{\gamma}\Psi_D(u,\Theta),\\
 \nu_{D,ei}^{\rm FT,0}
 &=\frac{1}{\tau_c}\frac{\gamma}{u^3}Z_{\rm eff},\\
 \nu_D^{\rm FT,0}
 &=\nu_{D,ee}^{\rm FT,0}+\nu_{D,ei}^{\rm FT,0}.
 \label{eq:physical_nud}
\end{align}
The finite-temperature relativistic functions are
\begin{align}
 \Psi_s
 &=\frac{
   \gamma^2\Psi_1-\Theta\Psi_0
   +(\Theta\gamma-1)u e^{-\gamma/\Theta}}
  {u^2K_2(1/\Theta)},
 \label{eq:psis_physical}\\
 \Psi_D
 &=\frac{
  (u^2\gamma^2+\Theta^2)\Psi_0
  +\Theta(2u^4-1)\Psi_1
  +\gamma\Theta[1+\Theta(2u^2-1)]u e^{-\gamma/\Theta}}
  {2\gamma u^3K_2(1/\Theta)},
 \label{eq:psid_physical}
\end{align}
with
\begin{equation}
 \Psi_0(u,\Theta)
 =\int_0^u\frac{e^{-\sqrt{1+s^2}/\Theta}}{\sqrt{1+s^2}}\,ds,
 \qquad
 \Psi_1(u,\Theta)
 =\int_0^u e^{-\sqrt{1+s^2}/\Theta}\,ds.
 \label{eq:psi_integrals_physical}
\end{equation}
Here $K_2$ is the modified Bessel function of the second kind.  Numerically stable
evaluation of Eqs.~\eqref{eq:psis_physical}--\eqref{eq:psi_integrals_physical}, including
the low-$u$ limit, belongs to the numerical model in Sec.~2 rather than to the physical
definition of the operator.

The thermal logarithm appearing in Eq.~\eqref{eq:physical_collision_time} is
\begin{equation}
 \ln\Lambda_0
 =14.9-\frac12\ln\left(\frac{n_e}{10^{20}\,{\rm m}^{-3}}\right)
 +\ln\left(\frac{T_e}{1\,{\rm keV}}\right).
 \label{eq:physical_thermal_log}
\end{equation}

\paragraph{Energy-dependent Coulomb-logarithm correction common to both base models.}
The thermal-to-relativistic matching of the Coulomb logarithms is a physical correction
layer, not part of the definition of either small-angle base kernel.  Following Hesslow
et al.~\cite{Hesslow2018Collision}, define
\begin{align}
 \ln\Lambda_{ee}(u)
 &=\ln\Lambda_0+\frac{1}{k}
 \ln\left[
  1+\left(\frac{2(\gamma-1)}{u_{Te}^2}\right)^{k/2}
 \right],\\
 \ln\Lambda_{ei}(u)
 &=\ln\Lambda_0+\frac{1}{k}
 \ln\left[1+\left(\frac{2u}{u_{Te}}\right)^k\right],
 \qquad
 u_{Te}=\sqrt{2\Theta},
 \qquad k=5,
 \label{eq:physical_matched_logs}
\end{align}
and the ratios
\begin{equation}
 R_{ee}(u)=\frac{\ln\Lambda_{ee}(u)}{\ln\Lambda_0},
 \qquad
 R_{ei}(u)=\frac{\ln\Lambda_{ei}(u)}{\ln\Lambda_0}.
 \label{eq:physical_log_ratios}
\end{equation}
For either base $b\in\{\mathrm{FT},\mathrm{match}\}$, decompose its thermal-log
complete-screening deflection frequency into electron--electron and electron--ion parts,
$\nu_D^{b,0}=\nu_{D,ee}^{b,0}+\nu_{D,ei}^{b,0}$.  The common Coulomb-logarithm layer is
then
\begin{align}
 \nu_s^{b,\rm CL}&=R_{ee}\,\nu_s^{b,0},\\
 \nu_\parallel^{b,\rm CL}&=\nu_\parallel^{b,0},\\
 \nu_D^{b,\rm CL}
 &=R_{ee}\,\nu_{D,ee}^{b,0}
  +R_{ei}\,\nu_{D,ei}^{b,0}.
 \label{eq:physical_common_log_overlay}
\end{align}
Thus the electron--electron logarithm corrects collisional drag and the
 electron--electron part of pitch-angle deflection, the electron--ion logarithm corrects
only the electron--ion deflection sector, and radial energy diffusion is unchanged.
For the finite-temperature base the split is explicit in
Eqs.~\eqref{eq:physical_nus}--\eqref{eq:physical_nud}.  For the matched base, the
$Z_{\rm eff}$ part of the deflection coefficient is the electron--ion sector and the
remaining deflection terms are the electron--electron sector, as made explicit in the
solver-normalized formulas in Sec.~2.

\paragraph{Partial-screening correction common to both base models.}
For a partially ionized ion species $a$ with atomic number $Z_a$, ion charge
$Z_{0a}$, density $n_a$, and number of bound electrons $N_a=Z_a-Z_{0a}$, define
\cite{Hesslow2017Screening,Hesslow2018Collision}
\begin{align}
 h(u)
 &=\sum_a\frac{n_a}{n_e}N_a
 \left\{
  \frac15\ln\left[
   1+\left(\frac{u\sqrt{\gamma-1}}{I_a/(m_ec^2)}\right)^5
  \right]-\beta^2
 \right\},
 \label{eq:physical_h}\\
 g(u)
 &=\sum_a\frac{n_a}{n_e}
 \left\{
  \frac23(Z_a^2-Z_{0a}^2)
       \ln\left[1+(u\bar a_a)^{3/2}\right]
  -\frac23N_a^2
       \frac{(u\bar a_a)^{3/2}}{1+(u\bar a_a)^{3/2}}
 \right\}.
 \label{eq:physical_g}
\end{align}
Here $I_a$ is the mean excitation energy and $\bar a_a$ is the dimensionless screening
length defined in Ref.~\cite{Hesslow2018Collision}.  The corresponding physical
frequency corrections are
\begin{align}
 \Delta\nu_s^{\rm PS}
 &=\frac{1}{\tau_c}\frac{\gamma^2}{u^3}
   \frac{h(u)}{\ln\Lambda_0},\\
 \Delta\nu_\parallel^{\rm PS}&=0,\\
 \Delta\nu_D^{\rm PS}
 &=\frac{1}{\tau_c}\frac{\gamma}{u^3}
   \frac{g(u)}{\ln\Lambda_0}.
 \label{eq:physical_partial_screening_overlay}
\end{align}
When state-resolved charge populations are supplied, each species sum in
Eqs.~\eqref{eq:physical_h}--\eqref{eq:physical_g} also sums over charge state.  For
charge state $q$, use
\begin{equation}
 Z_{0a}=q,\qquad N_{a,q}=Z_a-q,\qquad
 I_a=I_{a,q},\qquad \bar a_a=\bar a_{a,q},\qquad
 n_a=n_{a,q}.
 \label{eq:physical_resolved_screening_inputs}
\end{equation}
The code uses this state-resolved form in coupled mode.  Scalar $Z_{0a}$, $I_a$, and
$\bar a_a$ inputs remain available in standalone mode.
For neutral deuterium, the configured effective screening length is
$\bar a_{\mathrm{D^0}}=190$; the fully stripped state uses the zero sentinel because
its partial-screening correction vanishes.
For either base kernel, the partial-screening correction is applied after the common
energy-dependent Coulomb-logarithm layer.  Equivalently,
$\Delta C_F^{(P),\rm PS}=P\Delta\nu_s^{\rm PS}$,
$\Delta C_A^{(P),\rm PS}=0$, and the deflection coefficient receives
$\Delta\nu_D^{\rm PS}$.  The final physical frequencies are therefore
\begin{align}
 \nu_s^b&=\nu_s^{b,\rm CL}+\Delta\nu_s^{\rm PS},\\
 \nu_\parallel^b&=\nu_\parallel^{b,\rm CL}=\nu_\parallel^{b,0},\\
 \nu_D^b&=\nu_D^{b,\rm CL}+\Delta\nu_D^{\rm PS},
 \qquad b\in\{\mathrm{FT},\mathrm{match}\}.
 \label{eq:physical_final_collision_layers}
\end{align}
Thus neither the energy-dependent Coulomb-logarithm layer nor the partial-screening layer
selects or blends the base operator, and neither modifies radial energy diffusion.

The effective charge is
\begin{equation}
 Z_{\rm eff}=\frac{1}{n_e}\sum_a\sum_q n_{a,q}q^2,
 \qquad
 n_e=\sum_a\sum_q n_{a,q}q,
 \label{eq:physical_zeff_resolved}
\end{equation}
with the scalar-ion form recovered when each species has one nonzero charge state.
In the complete-screening limit $N_a=0$ for every species, the screening corrections
vanish identically, $h=g=0$, while the common energy-dependent Coulomb-logarithm layer
remains active.  Setting additionally $R_{ee}=R_{ei}=1$ recovers the corresponding
thermal-log base kernel exactly.


\paragraph{Fully-relativistic and asymptotically matched test-particle model.}
The second selectable thermal-log complete-screening base collision kernel is the fully-relativistic,
asymptotically matched test-particle operator.  Stahl et al. describe its construction as
an asymptotic matching of the usual non-relativistic test-particle Fokker--Planck
operator to the fully relativistic high-energy test-particle limit
\cite{Stahl2016}.  The resulting momentum-space coefficients retain relativistic
kinematics throughout the represented tail while assuming a non-relativistic thermal
bulk.  This is historically the test-particle operator used in CODE, but CODE is
provenance rather than the model name.  It is not the finite-temperature
Braams--Karney/Pike--Rose operator above, and it does not contain the field-particle part
of the linearized collision operator.

For provenance and comparison with Stahl et al., it is useful to state this model in
the normalized variables used in their CODE implementation.  Choose a reference thermal speed
\begin{equation}
 \widetilde v_e=\sqrt{\frac{2\widetilde T_e}{m_e}},\qquad
 \delta\equiv\frac{\widetilde v_e}{c},
\end{equation}
and define
\begin{equation}
 y\equiv\frac{P}{m_e\widetilde v_e}=\frac{u}{\delta},\qquad
 x\equiv\frac{y}{\gamma}=\frac{v}{\widetilde v_e},\qquad
 \overline v_e\equiv\frac{v_e}{\widetilde v_e},\qquad
 v_e=\sqrt{\frac{2T_e}{m_e}}.
 \label{eq:matched_variables}
\end{equation}
With time normalized to a reference electron--electron collision frequency
$\widetilde\nu_{ee}$, the fully-relativistic asymptotically matched test-particle operator is
\begin{equation}
 \boxed{
 \widehat{\mathcal C}_{\rm match}^{\rm tp}[G]
 =c_C\overline v_e^{\,3}y^{-2}
 \left\{
 \frac{\partial}{\partial y}
 \left[
  y^2\Psi
  \left(
   \frac{1}{x}\frac{\partial}{\partial y}
   +\frac{2}{\overline v_e^{\,2}}
  \right)G
 \right]
 +\frac{c_\xi}{2x}\frac{\partial}{\partial\xi}
 \left[(1-\xi^2)\frac{\partial G}{\partial\xi}\right]
 \right\},
 }
 \label{eq:matched_tp_operator}
\end{equation}
where a bar denotes normalization to the reference plasma state, in particular
$\overline\nu_{ee}=\nu_{ee}/\widetilde\nu_{ee}$, and
\begin{align}
 c_C&=\frac{3\sqrt{\pi}}{4}\,\overline\nu_{ee},\\
 c_\xi&=Z_{\rm eff}+\Phi-\Psi
       +\frac{\overline v_e^{\,2}\delta^4x^2}{2},
 \label{eq:matched_coefficients}
\end{align}
with
\begin{equation}
 \Phi=\operatorname{erf}\!\left(\frac{x}{\overline v_e}\right),\qquad
 \Psi=\frac{\overline v_e^{\,2}}{2x^2}
 \left[
  \Phi-\frac{x}{\overline v_e}
  \frac{d\Phi}{d(x/\overline v_e)}
 \right].
 \label{eq:matched_chandrasekhar}
\end{equation}
For a constant plasma chosen equal to the reference state,
$\overline v_e=\overline\nu_{ee}=1$.  The assumption $\delta\ll1$ is the
non-relativistic-bulk ordering of this model.  The role of
Eq.~\eqref{eq:matched_tp_operator} in the model hierarchy is specifically to
recover the standard non-relativistic test-particle coefficients at thermal and
suprathermal energies while matching to the relativistic test-particle asymptote at
large momentum \cite{Stahl2016}.  This provides a useful fully-relativistic asymptotically matched configuration for
benchmarks without replacing the finite-temperature relativistic baseline above.  For this matched base, the slowing-down coefficient is electron--electron, the
$Z_{\rm eff}$ contribution to pitch-angle deflection is electron--ion, and the remaining
pitch-angle-deflection terms are electron--electron.  Consequently the same common
Coulomb-logarithm layer in Eq.~\eqref{eq:physical_common_log_overlay} applies to this
base without introducing a matched-model-specific logarithm prescription.  When bound
electrons are present, the same Eq.~\eqref{eq:physical_partial_screening_overlay}
correction is then added; no separate screening prescription is introduced either.

\subsection{Synchrotron radiation-reaction operator}

The gyro-averaged synchrotron radiation-reaction force is represented by the
momentum-space characteristics
\begin{equation}
 \dot P_{\rm syn}
 =-\frac{\gamma P}{\tau_s}(1-\xi^2),
 \qquad
 \dot\xi_{\rm syn}
 =\frac{\xi}{\gamma\tau_s}(1-\xi^2),
 \label{eq:synch_characteristics}
\end{equation}
with
\begin{equation}
 \tau_s=\frac{6\pi\epsilon_0m_e^3c^3}{e^4B^2}.
 \label{eq:synch_time}
\end{equation}
The conservative operator is
\begin{equation}
 \mathcal R_{\rm syn}[F_e]
 =-\frac{1}{P^2}\frac{\partial}{\partial P}
   \left(P^2\dot P_{\rm syn}F_e\right)
  -\frac{\partial}{\partial\xi}
   \left(\dot\xi_{\rm syn}F_e\right).
 \label{eq:synch_operator_dimensional}
\end{equation}
This is the standard gyro-averaged synchrotron reaction used in homogeneous
runaway-electron kinetic models; see Refs.~\cite{Andersson2001}.

\subsection{Chiu--Harvey large-angle secondary-electron gain}

The production model uses a single optional large-angle secondary-electron gain model:
the finite-primary-energy, aligned-primary Chiu--Harvey source
\cite{Chiu1998,Harvey2000}.  The source is gain only: it adds the represented
secondary electron but does not include the compensating recoil/loss of the primary.
The source may be disabled by setting $\mathcal S_{\rm LA}=0$.

Introduce the dimensionless momenta
\begin{equation}
 u\equiv\frac{P}{m_ec},\qquad
 u'\equiv\frac{P'}{m_ec},\qquad
 \gamma=\sqrt{1+u^2},\qquad
 \gamma'=\sqrt{1+u'^2},
 \label{eq:ch_auxiliary_momenta}
\end{equation}
and the rescaled distribution
\begin{equation}
 \widehat F_e(u,\xi,t)\equiv(m_ec)^3F_e(P,\xi,t).
 \label{eq:ch_hat_definition}
\end{equation}
This is only a change of momentum coordinate; it is not the reference normalization of
Sec.~\ref{sec:normalization}.  Define the pitch-integrated primary spectrum
\begin{equation}
 \boxed{
 \mathcal F(u',t)=\int_{-1}^{1}\widehat F_e(u',\xi',t)\,d\xi'.
 }
 \label{eq:ch_primary_spectrum}
\end{equation}
The Chiu--Harvey approximation retains this evolving primary energy spectrum but uses a
field-aligned primary direction, $\xi'=-1$, in the relativistic electron--electron
knock-on kinematics.

The normalized M{\o}ller differential cross section underlying the Chiu--Harvey source is
\cite{Moller1932,Chiu1998,Harvey2000}
\begin{equation}
 \frac{d\overline\sigma_{\rm ee}(\gamma',\gamma)}{du}
 =2\pi
 \frac{\beta\,\gamma'^2}{(\gamma'-1)^3(\gamma'+1)}
 \left[
  x^2-3x+\left(\frac{\gamma'-1}{\gamma'}\right)^2(1+x)
 \right],
 \label{eq:ch_knockon_cross_section}
\end{equation}
where $\overline\sigma_{\rm ee}=\sigma_{\rm ee}/r_e^2$,
\begin{equation}
 \beta=\frac{u}{\gamma},\qquad
 x=\frac{(\gamma'-1)^2}{(\gamma-1)(\gamma'-\gamma)},\qquad
 r_e=\frac{e^2}{4\pi\epsilon_0m_ec^2}.
 \label{eq:ch_knockon_cross_section_aux}
\end{equation}
For the aligned primary, the secondary pitch is
\begin{equation}
 \xi_{\rm CH}(\gamma',\gamma)
 =-\left[
  \frac{(\gamma'+1)(\gamma-1)}{(\gamma'-1)(\gamma+1)}
 \right]^{1/2}.
 \label{eq:ch_secondary_pitch}
\end{equation}
The corresponding gain operator is
\begin{equation}
 \widehat S_{\rm CH}(u,\xi,t)
 =\frac{n_t c r_e^2}{u^2}
  \int_0^\infty du'\,u'^2\beta'
  \frac{d\overline\sigma_{\rm ee}(\gamma',\gamma)}{du}
  \delta\!\left[\xi-\xi_{\rm CH}(\gamma',\gamma)\right]
  \mathcal F(u',t),
 \label{eq:ch_integral_source}
\end{equation}
with $\beta'=u'/\gamma'$.  Thus no prescribed analytic primary spectrum is introduced:
$\mathcal F$ is obtained directly from the evolving 2P solution.

Equation~\eqref{eq:ch_integral_source} can be reduced analytically over $u'$.  The
required primary root is
\begin{equation}
 u'_0
 =-\frac{2u\xi}{1+\xi^2-\gamma(1-\xi^2)}
 =\frac{2u|\xi|}{1+\xi^2-\gamma(1-\xi^2)},
 \qquad \xi<0,
 \label{eq:ch_primary_root}
\end{equation}
with Jacobian
\begin{equation}
 \left|
 \frac{\partial(\xi-\xi_{\rm CH})}{\partial u'}
 \right|_{u'=u'_0}
 =\frac{|\xi|}{u'_0\gamma'_0},
 \qquad \gamma'_0=\sqrt{1+u_0'^2}.
 \label{eq:ch_delta_jacobian}
\end{equation}
Using $\beta'_0=u'_0/\gamma'_0$ gives
\begin{equation}
 \boxed{
 \widehat S_{\rm CH}(u,\xi,t)
 =\frac{n_t c r_e^2}{u^2}
  \frac{d\overline\sigma_{\rm ee}(u'_0,u)}{du}
  \frac{u_0'^4}{|\xi|}\,
  \mathcal F(u'_0,t),
 }
 \label{eq:ch_reduced_source}
\end{equation}
when the root satisfies the kinematic and domain constraints, and zero otherwise.

A large-angle cutoff $P_m$ avoids double counting scattering already represented
by the small-angle Fokker--Planck operator.  In the current prototype it is set equal
to the fixed kinetic lower boundary, $P_m=P_{\min}$.  With
\begin{equation}
 \epsilon_m=\sqrt{1+(P_m/m_ec)^2}-1,
\end{equation}
the retained secondary kinetic energy $\epsilon_s=\gamma-1$ for a primary with
$\epsilon_p=\gamma'-1$ satisfies
\begin{equation}
 \epsilon_s\in\left[\epsilon_m,\frac{\epsilon_p}{2}\right].
 \label{eq:ch_energy_range}
\end{equation}
The source is also zero if the root in Eq.~\eqref{eq:ch_primary_root} lies outside the
represented primary radial domain.

The dimensional large-angle term is therefore simply
\begin{equation}
 \boxed{
 \mathcal S_{\rm LA}[F_e](P,\xi,t)
 \in\left\{0,\;\frac{1}{(m_ec)^3}\widehat S_{\rm CH}(u,\xi,t)\right\}.
 }
 \label{eq:large_angle_physical_choices}
\end{equation}
The target density $n_t$ is explicit and denotes the population treated as cold
stationary target electrons for the knock-on source.  For the partially ionized
application of Ref.~, $n_t$ is the total free-plus-bound electron
density; in the standalone solver it is an explicit prescribed physical input.

\subsection{Prescribed plasma state, initial condition, and moments}

For standalone mode, the prescribed quantities
\begin{equation}
 \left\{T_e,n_e,E_\parallel,B,n_a,Z_a,Z_{0a},I_a,\bar a_a,n_t,P_{\min}\right\},
 \qquad P_m=P_{\min}
\end{equation}
are constant in time.  Consequently every coefficient in
Eq.~\eqref{eq:dimensional_kinetic} is time independent.  In coupled mode, the same
physical coefficient construction is reevaluated at each bulk stage.

The initial condition is an arbitrary non-negative gyrotropic distribution
\begin{equation}
 F_e(P,\xi,0)=F_{e0}(P,\xi).
 \label{eq:dimensional_initial_condition}
\end{equation}
Useful physical moments include the represented density in
Eq.~\eqref{eq:dimensional_density} and the parallel current density
\begin{equation}
 j_\parallel(t)
 =-e\,2\pi\int_0^\infty P^2\,dP\int_{-1}^{1}
  v(P)\,\xi\,F_e(P,\xi,t)\,d\xi.
 \label{eq:dimensional_current}
\end{equation}
Additional moments may be formed from the same distribution without modifying the
kinetic equation.  The total kinetic current in Eq.~\eqref{eq:dimensional_current}
is a diagnostic moment and is not, by itself, a threshold-defined runaway-electron
current.

For a user-selected diagnostic threshold $\kappa$, a thresholded runaway current
would instead be defined by
\begin{equation}
 j_{\rm RE}^{(\kappa)}(t)
 =-e\,2\pi\int_0^\infty P^2\,dP\int_{-1}^{1}
 v(P)\,\xi\,F_e(P,\xi,t)\,
 \mathbf 1_{\{\gamma(P)-1>\kappa T_e/(m_ec^2)\}}\,d\xi,
 \qquad
 I_{\rm RE}^{(\kappa)}=A j_{\rm RE}^{(\kappa)}.
 \label{eq:threshold_runaway_current}
\end{equation}
The full moment $j_\parallel$ remains available as a diagnostic.  The provisional
CPU coupling uses the thresholded non-thermal moment in
Eq.~\eqref{eq:threshold_runaway_current}; the precise current closure remains an
open model decision and must be qualified against the full-moment alternative.  The
Chiu--Harvey cutoff is set equal to the kinetic lower boundary,
$P_m=P_{\min}$, in the current prototype.

The implemented Maxwell--J{\"u}ttner initial condition uses the dimensionless momentum
$p=P/(m_ec)$ and temperature $\Theta_i=T_i/(m_ec^2)$.  Its unnormalised isotropic shape is
\begin{equation}
 F_{\rm MJ}(p,\xi;T_i)\propto
 \exp\left[-\frac{\sqrt{1+p^2}-1}{\Theta_i}\right],
 \qquad -1\leq\xi\leq1,
 \label{eq:maxwell_juttner_initial}
\end{equation}
and the finite-volume projection normalises the represented density to the configured
$n_{\rm init}$ (or $n_e$).  An optional seed is mixed by particle content; it does not
increase the configured total density.

\subsection{Normalization used by the numerical model}
\label{sec:normalization}

Only after the dimensional physical model has been defined do we introduce the
normalization used in Sec.~2.  The local flux form below is common to both specified
test-particle small-angle configurations; Sec.~2 supplies the model-specific values of
$C_F$, $C_A$, and $\nu_D$.  The large-angle term is selected independently as either none or Chiu--Harvey.  Choose a fixed reference density $n_{\rm ref}$
and reference Coulomb logarithm $\ln\Lambda_{\rm ref}$ and define
\begin{equation}
 \tau_{\rm ref}
 =\frac{4\pi\epsilon_0^2m_e^2c^3}
 {n_{\rm ref}e^4\ln\Lambda_{\rm ref}},
 \qquad
 E_{\rm ref}=\frac{m_ec}{e\tau_{\rm ref}}.
 \label{eq:reference_scales}
\end{equation}
The dimensionless numerical variables are
\begin{equation}
 p=\frac{P}{m_ec},
 \qquad
 \tau=\frac{t}{\tau_{\rm ref}},
 \qquad
 f(p,\xi,\tau)=\frac{(m_ec)^3}{n_{\rm ref}}F_e(P,\xi,t),
 \qquad
 \bar E=\frac{E_\parallel}{E_{\rm ref}}.
 \label{eq:numerical_normalization}
\end{equation}
The normalization is fixed for the calculation and does not alter the dimensional
physics.  In particular,
\begin{equation}
 \frac{n_F}{n_{\rm ref}}
 =2\pi\int_0^\infty p^2\,dp\int_{-1}^{1}f(p,\xi,\tau)\,d\xi.
 \label{eq:normalized_density}
\end{equation}
Define normalized collision coefficients
\begin{equation}
 C_F=p\tau_{\rm ref}\nu_s,
 \qquad
 C_A=\frac{p^2}{2}\tau_{\rm ref}\nu_\parallel,
 \qquad
 \nu_D^{\ast}=\tau_{\rm ref}\nu_D,
 \qquad
 \alpha=\frac{\tau_{\rm ref}}{\tau_s}.
 \label{eq:normalized_coefficients}
\end{equation}
In the remainder of the document the star is dropped and $\nu_D$ denotes the
normalized deflection frequency.

With these definitions the normalized forward equation used by the numerical model is
\begin{equation}
 \frac{\partial f}{\partial\tau}
 =-\frac{1}{p^2}\frac{\partial}{\partial p}\left(p^2J_p\right)
  -\frac{\partial J_\xi}{\partial\xi}
  +S_{\rm LA}[f],
 \label{eq:normalized_forward}
\end{equation}
with local fluxes
\begin{align}
 J_p
 &=A_p f-C_A\frac{\partial f}{\partial p},
 \label{eq:normalized_Jp}\\
 J_\xi
 &=A_\xi f
   -\frac{\nu_D}{2}(1-\xi^2)\frac{\partial f}{\partial\xi},
 \label{eq:normalized_Jxi}
\end{align}
and
\begin{align}
 A_p
 &=-\bar E\,\xi
   -\alpha\gamma p(1-\xi^2)
   -C_F,
 \label{eq:normalized_Ap}\\
 A_\xi
 &=(1-\xi^2)
   \left(-\frac{\bar E}{p}+\alpha\frac{\xi}{\gamma}\right).
 \label{eq:normalized_Axi}
\end{align}
The normalized knock-on event-rate scale follows directly from
Eq.~\eqref{eq:ch_integral_source}.  Factoring the $2\pi$ contained in the normalized
differential cross section in Eq.~\eqref{eq:ch_knockon_cross_section},
\begin{equation}
 \tau_{\rm ref}\,2\pi n_t c r_e^2
 =\frac{n_t/n_{\rm ref}}{2\ln\Lambda_{\rm ref}}.
 \label{eq:normalized_knockon_scale}
\end{equation}
This normalization is used by the Chiu--Harvey source.  For constant prescribed plasma
parameters, all selected local coefficients and the Chiu--Harvey map are constant in
time.  The coupled driver uses the same normalization while updating stage-dependent
plasma coefficients.

\section{Numerical model}
%===============================================================================

The standalone numerical problem is the constant-coefficient normalized forward equation
\eqref{eq:normalized_forward}.  The current CPU prototype uses the fully-relativistic
asymptotically matched thermal-log test-particle base coefficient set.  The common
energy-dependent Coulomb-logarithm layer is then applied to the electron--electron and
electron--ion sectors, and bound electrons activate the common partial-screening
drag/deflection layer defined in Sec.~1.  Fully stripped species therefore remove only
the partial-screening layer; the energy-dependent Coulomb logarithms remain part of the
physical collision coefficients.  The finite-temperature base is retained in this
document as a reference formulation for a future implementation.
The large-angle gain is either disabled or Chiu--Harvey.  Every configuration uses the same physical
finite-volume state in the phase-space measure $2\pi p^2\,dp\,d\xi$ and is advanced as
part of one fully implicit linear kinetic operator.  For fixed physical parameters,
configuration, mesh, and timestep, the standalone semidiscrete operator is linear and
time independent.  The coupled driver retains linear kinetic stages but updates
coefficients from host-side bulk stage states.

\subsection{Finite-volume grid, representative points, and discrete state}

The computational domain is
\begin{equation}
 p_{\min}\le p\le p_{\max},\qquad -1\le\xi\le1,
\end{equation}
and is discretized with a tensor-product mesh.  The current CPU implementation uses
a hybrid radial mesh and a hybrid pitch mesh.  Thus $p_{\min}>0$ is required by the
present logarithmic low-momentum segment; the current production configuration uses
$p_{\min}=100\,{\rm eV}/(m_ec^2)$.  More generally, for the fixed nonthermal factor
$\kappa=10$, the current cutoff is
\begin{equation}
 p_{\min}=\sqrt{\left(1+\frac{\kappa T_e}{m_ec^2}\right)^2-1}.
 \label{eq:current_nonthermal_cutoff}
\end{equation}
The radial faces use a logarithmic segment from $p_{\min}$ to the fixed transition
$p_\star=1$, followed by a uniform-in-$p$ segment from $p_\star$ to $p_{\max}$.
The number of low-momentum cells is allocated in proportion to logarithmic range,
\begin{equation}
 N_{\rm low}=\operatorname{round}\left[
 N_p\frac{\ln(p_\star/p_{\min})}{\ln(p_{\max}/p_{\min})}
 \right],
 \qquad N_{\rm high}=N_p-N_{\rm low}.
 \label{eq:hybrid_radial_allocation}
\end{equation}
The radial faces and cell centers in each segment are
\begin{equation}
 p_{i+1/2}=p_{\min}\left(\frac{p_\star}{p_{\min}}\right)^{i/N_{\rm low}}
 \quad (0\le i\le N_{\rm low}),
 \qquad
 p_{N_{\rm low}+m+1/2}=p_\star+
 \frac{p_{\max}-p_\star}{N_{\rm high}}m
 \quad (0\le m\le N_{\rm high}),
 \label{eq:geometric_p_grid}
\end{equation}
and $p_i=\sqrt{p_{i-1/2}p_{i+1/2}}$ is used as the representative point in both
segments.
The distances used by the radial flux stencil are
\begin{equation}
 d_{p,1/2}=p_1-p_{1/2},\qquad
 d_{p,i+1/2}=p_{i+1}-p_i\;(1\le i<N_p),\qquad
 d_{p,N_p+1/2}=p_{N_p+1/2}-p_{N_p}.
 \label{eq:geometric_p_distances}
\end{equation}
For pitch angle, the negative half uses uniformly spaced $\theta$ and the positive
half remains uniformly spaced in $\xi$.  With
$N_-=\lfloor N_\xi/2\rfloor$ and $N_+=N_\xi-N_-$,
\begin{equation}
 \xi_{j+1/2}=\cos\theta_{N_- -j},
 \quad \theta_k=\frac{\pi}{2}+\frac{k\pi}{2N_-},
 \quad 0\le j<N_-,
 \qquad
 \xi_{N_-+m+1/2}=\frac{m}{N_+},
 \quad 0\le m\le N_+,
 \label{eq:current_xi_grid}
\end{equation}
The pitch cell widths and center-to-face distances are
\begin{equation}
 \Delta\xi_j=\xi_{j+1/2}-\xi_{j-1/2},\qquad
 d_{\xi,1/2}=\xi_1-\xi_{1/2},\qquad
 d_{\xi,j+1/2}=\xi_{j+1}-\xi_j,
 \qquad d_{\xi,N_\xi+1/2}=\xi_{N_\xi+1/2}-\xi_{N_\xi}.
 \label{eq:current_xi_distances}
\end{equation}
The radial mapping is a living numerical choice: it will be replaced or retuned after
momentum-resolution, boundary, and $p_{\max}$ convergence studies.  Resolution is
obtained by increasing $N_p$ and $N_\xi$; the local radial spacing must resolve both
the low-energy collision layer and the relativistic drift/radiation region.

The primary finite-volume unknown is the exact normalized particle content of cell
$(i,j)$,
\begin{equation}
 N_{ij}(\tau)
 =2\pi\int_{p_{i-1/2}}^{p_{i+1/2}}p^2\,dp
       \int_{\xi_{j-1/2}}^{\xi_{j+1/2}}f(p,\xi,\tau)\,d\xi.
 \label{eq:cellcontent}
\end{equation}
The exact cell phase volume and piecewise-constant reconstruction are
\begin{align}
 V_{ij}
 &=\frac{2\pi}{3}
   \left(p_{i+1/2}^3-p_{i-1/2}^3\right)
   \Delta\xi_j,
 \label{eq:cellvolume}\\
 f_{ij}&=\frac{N_{ij}}{V_{ij}}.
 \label{eq:cell_reconstruction}
\end{align}
The PETSc implementation stores $N_{ij}$ as the primary kinetic state and reconstructs
$f_{ij}$ only when evaluating local fluxes, coefficients, currents, or diagnostics.
The current logarithmic grid has $p_{\min}>0$, so the first cell terminates on the
prescribed inner boundary.  A future origin-touching grid with $p_{\min}=0$ would impose
regularity through the vanishing radial surface area and would require no coefficient at
$p=0$ itself.

The continuous initial condition is projected into this same finite-volume measure,
\begin{equation}
 N_{ij}^{0}
 =2\pi\int_{p_{i-1/2}}^{p_{i+1/2}}p^2\,dp
       \int_{\xi_{j-1/2}}^{\xi_{j+1/2}}f_0(p,\xi)\,d\xi.
 \label{eq:initial_projection}
\end{equation}
The projection is evaluated analytically when possible and otherwise by converged
Gauss--Legendre quadrature; point sampling of $f_0$ at the cell center is not used as
a substitute for Eq.~\eqref{eq:initial_projection}.

\subsection{Numerical evaluation of the fixed physical coefficients}

All physical coefficients are time independent for the present problem and are therefore
evaluated once before the kinetic operator is assembled.  The current CPU implementation
reduces the fully-relativistic asymptotically matched base and its common correction
layers to the three coefficient arrays $C_F(p)$, $C_A(p)$, and $\nu_D(p)$ consumed by
the finite-volume fluxes below.  The coefficient pipeline is
\begin{equation}
 \begin{aligned}
 \text{thermal-log base}
 &\;\longrightarrow\;\text{common energy-dependent Coulomb logs}\\
 &\;\longrightarrow\;\text{common partial screening}
 \;\longrightarrow\;\{C_F,C_A,\nu_D\}.
 \end{aligned}
 \label{eq:collision_layer_pipeline}
\end{equation}
Selecting a collision model therefore changes only the base coefficient generation, not
the two common correction layers or the finite-volume stencil.  Define once for both
bases
\begin{equation}
 \chi_c=\frac{n_e}{n_{\rm ref}}
        \frac{\ln\Lambda_0}{\ln\Lambda_{\rm ref}}.
 \label{eq:common_collision_scaling}
\end{equation}

\paragraph{Finite-temperature relativistic coefficient evaluation (future reference).}
The finite-temperature collision functions in
Eqs.~\eqref{eq:psis_physical}--\eqref{eq:psi_integrals_physical} must be evaluated in a
form that remains accurate both at low temperature and as $p\rightarrow0$.  The
numerical treatment below is an alternative evaluation of the same functions defined in
Sec.~1; it does not introduce a different collision model.

\subparagraph{Exponentially scaled rapidity representation.}
With
\begin{equation}
 p=\sinh y,\qquad \gamma=\cosh y,
 \label{eq:rapidity_substitution}
\end{equation}
define
\begin{align}
 \widetilde\Psi_0(p,\Theta)
 &\equiv e^{1/\Theta}\Psi_0(p,\Theta)
 =\int_0^{\operatorname{asinh}p}
   e^{-[\cosh y-1]/\Theta}\,dy,
 \label{eq:psi0_scaled}\\
 \widetilde\Psi_1(p,\Theta)
 &\equiv e^{1/\Theta}\Psi_1(p,\Theta)
 =\int_0^{\operatorname{asinh}p}
   e^{-[\cosh y-1]/\Theta}\cosh y\,dy,
 \label{eq:psi1_scaled}\\
 \widetilde K_2(\Theta)
 &\equiv e^{1/\Theta}K_2(1/\Theta)
 =\int_0^\infty
   e^{-[\cosh y-1]/\Theta}\cosh(2y)\,dy.
 \label{eq:k2_scaled}
\end{align}
The scaled collision functions are therefore evaluated directly as
\begin{align}
 \Psi_s^{\rm dir}(p,\Theta)
 &=\frac{
   \gamma^2\widetilde\Psi_1-\Theta\widetilde\Psi_0
   +(\Theta\gamma-1)p e^{-(\gamma-1)/\Theta}}
  {p^2\widetilde K_2},
 \label{eq:psis_scaled_direct}\\
 \Psi_D^{\rm dir}(p,\Theta)
 &=\frac{
  (p^2\gamma^2+\Theta^2)\widetilde\Psi_0
  +\Theta(2p^4-1)\widetilde\Psi_1
  +\gamma\Theta[1+\Theta(2p^2-1)]p e^{-(\gamma-1)/\Theta}}
  {2\gamma p^3\widetilde K_2}.
 \label{eq:psid_scaled_direct}
\end{align}
Thus neither $e^{-1/\Theta}$ nor $K_2(1/\Theta)$ is formed separately.  The finite
integrals in Eqs.~\eqref{eq:psi0_scaled} and \eqref{eq:psi1_scaled} are evaluated by
Gauss--Legendre quadrature in $y$; the collision quadrature order $q_{\rm coll}$ is a
numerical convergence parameter.  The same converged value of $\widetilde K_2$ is used
in both the direct and low-$p$ branches.

\subparagraph{Low-$p$ series.}
Although Eqs.~\eqref{eq:psis_scaled_direct} and \eqref{eq:psid_scaled_direct} are free of
thermal underflow, their numerators contain cancellations as $p\rightarrow0$.  The
functions are therefore evaluated near the origin by their Taylor series obtained
directly from the same equations.  Write
\begin{align}
 \Psi_s^{\rm ser}(p,\Theta)
 &=\frac{p}{\widetilde K_2}
   \sum_{m=0}^{6}a_{2m+1}(\Theta)p^{2m},
 \label{eq:psis_lowp_series}\\
 \Psi_D^{\rm ser}(p,\Theta)
 &=\frac{1}{\widetilde K_2}
   \sum_{m=0}^{6}b_{2m}(\Theta)p^{2m}.
 \label{eq:psid_lowp_series}
\end{align}
The retained coefficients, corresponding to terms through $p^{13}$ in $\Psi_s$ and
$p^{12}$ in $\Psi_D$, are
\begin{align}
 a_1&=\frac{2\Theta^2+2\Theta+1}{3\Theta},\\
 a_3&=-\frac{6\Theta^3+6\Theta^2+5\Theta+3}{30\Theta^2},\\
 a_5&=\frac{30\Theta^4+30\Theta^3+27\Theta^2+17\Theta+5}{280\Theta^3},\\
 a_7&=-\frac{210\Theta^5+210\Theta^4+195\Theta^3+125\Theta^2+44\Theta+7}
 {3024\Theta^4},\\
 a_9&=\frac{1890\Theta^6+1890\Theta^5+1785\Theta^4+1155\Theta^3
 +435\Theta^2+92\Theta+9}{38016\Theta^5},\\
 a_{11}&=-\frac{20790\Theta^7+20790\Theta^6+19845\Theta^5+12915\Theta^4
 +5040\Theta^3+1197\Theta^2+167\Theta+11}{549120\Theta^6},\\
 a_{13}&=\frac{270270\Theta^8+270270\Theta^7+259875\Theta^6+169785\Theta^5
 +67725\Theta^4+17136\Theta^3+2786\Theta^2+275\Theta+13}
 {8985600\Theta^7},
 \label{eq:psis_lowp_coeffs}
\end{align}
and
\begin{align}
 b_0&=\frac{8\Theta^2+5\Theta+4}{12},\\
 b_2&=\frac{64\Theta^3+79\Theta^2+30\Theta-8}{240\Theta},\\
 b_4&=-\frac{768\Theta^4+873\Theta^3+274\Theta^2+53\Theta-12}
 {3360\Theta^2},\\
 b_6&=\frac{49152\Theta^5+53877\Theta^4+15594\Theta^3+3830\Theta^2
 +572\Theta-80}{241920\Theta^3},\\
 b_8&=-\frac{983040\Theta^6+1055805\Theta^5+292986\Theta^4+75031\Theta^3
 +16876\Theta^2+1545\Theta-140}{5322240\Theta^4},\\
 b_{10}&=\frac{47185920\Theta^7+50023755\Theta^6+13539990\Theta^5
 +3487020\Theta^4+866196\Theta^3+136817\Theta^2+8050\Theta-504}
 {276756480\Theta^5},\\
 b_{12}&=-\frac{
 \begin{aligned}
  &440401920\Theta^8+462699195\Theta^7+123237750\Theta^6
   +31617235\Theta^5+8090500\Theta^4\\
  &\qquad{}+1515314\Theta^3+167034\Theta^2+6755\Theta-308
 \end{aligned}
 }{2767564800\Theta^6}.
 \label{eq:psid_lowp_coeffs}
\end{align}
For numerical evaluation the two polynomials are evaluated in $z=p^2$ using Horner's
rule.

The baseline branch switch is fixed as
\begin{equation}
 p_{\rm sw}(\Theta)=0.1\min\!\left(1,\sqrt{\Theta}\right),
 \label{eq:lowp_switch}
\end{equation}
with
\begin{equation}
 (\Psi_s,\Psi_D)=
 \begin{cases}
  (\Psi_s^{\rm ser},\Psi_D^{\rm ser}), & p\le p_{\rm sw}(\Theta),\\
  (\Psi_s^{\rm dir},\Psi_D^{\rm dir}), & p>p_{\rm sw}(\Theta).
 \end{cases}
 \label{eq:collision_function_branch}
\end{equation}
This switch places the transition in a region where both $p$ and the thermal-scaled
momentum $p/\sqrt{\Theta}$ are small.  It is fixed for the finite-temperature configuration rather than
being treated as a user-adjustable modeling parameter.  The implementation must recover
\begin{equation}
 \Psi_s=O(p),\qquad \Psi_D=O(1),\qquad p\rightarrow0,
 \label{eq:low_p_limits}
\end{equation}
and the direct and series branches must agree in the neighborhood of
Eq.~\eqref{eq:lowp_switch} to the accuracy required of the coefficient evaluation.

The corresponding normalized finite-temperature thermal-log base coefficients are
\begin{align}
 C_F^{\rm FT,0}(p)
 &=\chi_c\,\Psi_s(p,\Theta),
 \label{eq:ft_numeric_cf_base}\\
 C_A^{\rm FT,0}(p)
 &=\chi_c\,\frac{\gamma\Theta}{p}\Psi_s(p,\Theta),
 \label{eq:ft_numeric_ca_base}\\
 \nu_{D,ee}^{\rm FT,0}(p)
 &=\chi_c\,\frac{2\Psi_D(p,\Theta)}{p^2},\\
 \nu_{D,ei}^{\rm FT,0}(p)
 &=\chi_c\,\frac{\gamma Z_{\rm eff}}{p^3},\\
 \nu_D^{\rm FT,0}(p)
 &=\nu_{D,ee}^{\rm FT,0}(p)+\nu_{D,ei}^{\rm FT,0}(p).
 \label{eq:ft_numeric_nud_base}
\end{align}
These arrays contain the base-kernel dependence only; the common Coulomb-logarithm and
partial-screening layers are applied below.

\paragraph{Fully-relativistic asymptotically matched coefficient evaluation (current CPU operator).}
For the current CPU configuration the reference plasma is chosen equal to the prescribed
constant plasma state, so that $\overline v_e=\overline\nu_{ee}=1$ in
Eqs.~\eqref{eq:matched_variables}--\eqref{eq:matched_coefficients}.  Define
\begin{equation}
 \delta=\sqrt{\frac{2T_e}{m_ec^2}},\qquad
 x(p)=\frac{p}{\delta\gamma}.
 \label{eq:matched_numeric_scalings}
\end{equation}
Transforming the fully-relativistic asymptotically matched operator from
$y=p/\delta$ and its historical collision-time normalization to the common $(p,\tau)$
variables of Sec.~\ref{sec:normalization} gives
exactly the same conservative radial/pitch form used by the finite-volume solver.  Using
the reference collision-frequency definition of Stahl et al.,
\begin{equation}
 \widetilde\nu_{ee}
 =\frac{16\sqrt{\pi}\,n_e e^4\ln\Lambda_0}
 {3m_e^2v_e^3(4\pi\epsilon_0)^2},
 \qquad
 (\tau_{\rm ref}\widetilde\nu_{ee})\frac{3\sqrt{\pi}}{4}
 =\chi_c\delta^{-3}.
 \label{eq:matched_numeric_frequency_conversion}
\end{equation}
The resulting thermal-log common-normalization base coefficients are
\begin{align}
 C_F^{\rm match,0}(p)
 &=\chi_c\,\frac{2\Psi(x)}{\delta^2},
 \label{eq:matched_numeric_cf}\\
 C_A^{\rm match,0}(p)
 &=\chi_c\,\frac{\gamma}{p}\Psi(x),
 \label{eq:matched_numeric_ca}\\
 \nu_{D,ee}^{\rm match,0}(p)
 &=\chi_c\,\frac{\gamma}{p^3}
 \left[
  \Phi(x)-\Psi(x)+\frac{\delta^2p^2}{2\gamma^2}
 \right],\\
 \nu_{D,ei}^{\rm match,0}(p)
 &=\chi_c\,\frac{\gamma}{p^3}Z_{\rm eff},\\
 \nu_D^{\rm match,0}(p)
 &=\nu_{D,ee}^{\rm match,0}(p)+\nu_{D,ei}^{\rm match,0}(p).
 \label{eq:matched_numeric_nud}
\end{align}
Here
\begin{equation}
 \Phi(x)=\operatorname{erf}(x),\qquad
 \Psi(x)=\frac{\Phi(x)-2x e^{-x^2}/\sqrt{\pi}}{2x^2}.
 \label{eq:matched_numeric_phi_psi}
\end{equation}

The CPU prototype uses these matched expressions with its local collision-time
normalization, equivalent to $n_{\rm ref}=n_e$, $\ln\Lambda_{\rm ref}=\ln\Lambda_0$,
and hence $\chi_c=1$.  Including the implemented common logarithm and partial-screening
layers, the actual arrays passed to the finite-volume operator are
\begin{align}
 C_F(p)&=R_{ee}(p)\frac{2\Psi(x)}{\delta^2}
       +\frac{\gamma^2h(p)}{p^2\ln\Lambda_0},\\
 C_A(p)&=\frac{\gamma}{p}\Psi(x),\\
 \nu_D(p)&=\frac{\gamma}{p^3}\left[
 R_{ee}(p)\left(\Phi(x)-\Psi(x)+\frac{\Theta p^2}{\gamma^2}\right)
 +Z_{\rm eff}R_{ei}(p)+\frac{g(p)}{\ln\Lambda_0}\right].
 \label{eq:cpu_matched_final_coefficients}
\end{align}
These are evaluated directly at the geometric-grid faces for radial fluxes and at cell
centers for pitch-angle diffusion.  The finite-temperature rapidity quadrature and
low-$p$ series are not used by the current CPU prototype.

\paragraph{Common energy-dependent Coulomb-logarithm layer in solver normalization.}
At every radial evaluation point, compute the same physical ratios as in
Eq.~\eqref{eq:physical_log_ratios}, with $u=p$,
\begin{equation}
 R_{ee}(p)=\frac{\ln\Lambda_{ee}(p)}{\ln\Lambda_0},
 \qquad
 R_{ei}(p)=\frac{\ln\Lambda_{ei}(p)}{\ln\Lambda_0}.
 \label{eq:numeric_log_ratios}
\end{equation}
For either $b\in\{\mathrm{FT},\mathrm{match}\}$ the Coulomb-logarithm-corrected
complete-screening arrays are
\begin{align}
 C_F^{b,\rm CL}(p)
 &=R_{ee}(p)\,C_F^{b,0}(p),
 \label{eq:common_log_numeric_cf}\\
 C_A^{b,\rm CL}(p)
 &=C_A^{b,0}(p),
 \label{eq:common_log_numeric_ca}\\
 \nu_D^{b,\rm CL}(p)
 &=R_{ee}(p)\,\nu_{D,ee}^{b,0}(p)
  +R_{ei}(p)\,\nu_{D,ei}^{b,0}(p).
 \label{eq:common_log_numeric_nud}
\end{align}
This step is identical for the two base choices.  In particular, changing the small-angle
base model does not disable, replace, or redefine either energy-dependent Coulomb
logarithm, and $C_A$ is unchanged by this layer.

\paragraph{Common partial-screening overlay in solver normalization.}
The same screening layer is applied after the common energy-dependent Coulomb-logarithm
correction for either base coefficient set.  In the common normalized variables its
contribution is
\begin{align}
 \Delta C_F^{\rm PS}(p)
 &=\chi_c\,\frac{\gamma^2}{p^2}\frac{h(p)}{\ln\Lambda_0},
 \label{eq:screening_numeric_cf}\\
 \Delta C_A^{\rm PS}(p)&=0,
 \label{eq:screening_numeric_ca}\\
 \Delta\nu_D^{\rm PS}(p)
 &=\chi_c\,\frac{\gamma}{p^3}\frac{g(p)}{\ln\Lambda_0}.
 \label{eq:screening_numeric_nud}
\end{align}
Thus, for either $b\in\{\mathrm{FT},\mathrm{match}\}$,
\begin{equation}
 \boxed{
 C_F=C_F^{b,\rm CL}+\Delta C_F^{\rm PS},\qquad
 C_A=C_A^{b,\rm CL}=C_A^{b,0},\qquad
 \nu_D=\nu_D^{b,\rm CL}+\Delta\nu_D^{\rm PS}.
 }
 \label{eq:common_screened_coefficients}
\end{equation}
The two base evaluations above therefore feed the finite-volume assembly through exactly
the same Coulomb-logarithm and screening semantics.  When every ion is fully stripped,
$h=g=0$ and Eq.~\eqref{eq:common_screened_coefficients} reduces to the selected base with
the common energy-dependent Coulomb-logarithm correction still active.  If one also sets
$R_{ee}=R_{ei}=1$, the corresponding thermal-log base is recovered exactly.  Disabling
radial energy diffusion remains a separate explicit model switch and is applied after
coefficient construction.
Equations~\eqref{eq:matched_numeric_cf}--\eqref{eq:matched_numeric_nud} are a coordinate
transformation of the Section-1 fully-relativistic asymptotically matched thermal-log base,
not a new interpolation or fit.  The common logarithm corrections in
Eqs.~\eqref{eq:common_log_numeric_cf}--\eqref{eq:common_log_numeric_nud} and the common
partial-screening terms in Eqs.~\eqref{eq:screening_numeric_cf}--\eqref{eq:screening_numeric_nud}
are then applied independently of the selected base.  The local finite-volume assembly
receives the same three final coefficient arrays in either case.  The production default
keeps radial energy diffusion.  A benchmark configuration may deliberately set $C_A=0$
after coefficient evaluation; this is an explicit model switch and is never activated
implicitly by mesh, temperature, screening, or large-angle settings.

No rapidity quadrature or modified Bessel function is required for the current CPU
configuration.  The code evaluates the Chandrasekhar functions directly from
Eq.~\eqref{eq:matched_chandrasekhar} at the geometric-grid faces and centers, then
applies the common logarithm and partial-screening overlays above.
For the current shifted domain all coefficient faces lie at finite momentum, so there is
no separate origin quadrature or singular coefficient evaluation to perform.  An
origin-touching grid remains a future grid variant and would impose its inner flux
directly.

For qualification, the transformed thermal-log base coefficients must reproduce the
Section-1 fully-relativistic asymptotically matched operator pointwise and recover its
intended asymptotic limits.  In particular, for $p\gg1$ with $\delta\ll1$,
\begin{equation}
 C_F^{\rm match,0}\rightarrow\chi_c,\qquad
 C_A^{\rm match,0}\rightarrow\frac{\chi_c\delta^2}{2},\qquad
 \nu_D^{\rm match,0}\rightarrow
 \chi_c\frac{1+Z_{\rm eff}}{p^2}
 \quad\text{to leading order in }\delta.
 \label{eq:matched_relativistic_limits}
\end{equation}
After the common Coulomb-logarithm layer, the corresponding complete-screening leading
behavior is instead
\begin{equation}
 C_F^{\rm match,CL}\rightarrow\chi_c R_{ee}(p),\qquad
 \nu_D^{\rm match,CL}\rightarrow
 \frac{\chi_c}{p^2}\left[R_{ee}(p)+Z_{\rm eff}R_{ei}(p)\right],
 \label{eq:matched_relativistic_limits_logs}
\end{equation}
while $C_A$ retains the base limit.  Partial screening adds the independent $h$ and $g$
terms afterward.  The non-relativistic thermal/suprathermal limit of the base is checked
directly against Eq.~\eqref{eq:matched_tp_operator} rather than by introducing a second
numerical formula.

\subsection{Conservative local flux discretization}

Integrating the differential part of Eq.~\eqref{eq:normalized_forward} over a cell gives
\begin{equation}
 \frac{dN_{ij}}{d\tau}\bigg|_{\rm local}
 =-\left(F^p_{i+1/2,j}-F^p_{i-1/2,j}\right)
  -\left(F^\xi_{i,j+1/2}-F^\xi_{i,j-1/2}\right).
 \label{eq:fv_balance}
\end{equation}
Each interior face flux is formed once and enters the two adjacent cells with opposite
signs.  Interior conservation is therefore algebraic rather than approximate.

The local drift--diffusion fluxes use the exponentially fitted Bernoulli form of the
Chang--Cooper discretization.  When the diffusion coefficient vanishes, the implementation
reduces to first-order upwinding.  This deliberately simple conservative choice is the
baseline numerical scheme for the present solver.  Higher-order schemes are outside the
scope of the present model and may be considered only after this formulation is fully
qualified.

Define the numerically stable Bernoulli function
\begin{equation}
 \mathcal B(z)=\frac{z}{e^z-1},\qquad \mathcal B(0)=1.
 \label{eq:bernoulli_function}
\end{equation}

For an interior radial face, define
\begin{equation}
 A^p_{i+1/2,j}=A_p(p_{i+1/2},\xi_j),\qquad
 C_{A,i+1/2}=C_A^{({\rm sel})}(p_{i+1/2}),
\end{equation}
and the face Peclet argument
\begin{equation}
 w^p_{i+1/2,j}=\frac{A^p_{i+1/2,j}d_{p,i+1/2}}{C_{A,i+1/2}}.
 \label{eq:radial_peclet}
\end{equation}
For $C_{A,i+1/2}>0$, the integrated radial face flux is
\begin{equation}
 F^p_{i+1/2,j}
 =2\pi p_{i+1/2}^2\Delta\xi
 \left[
   \frac{C_{A,i+1/2}}{d_{p,i+1/2}}
 \left(
   \mathcal B(-w^p_{i+1/2,j})f_{ij}
  -\mathcal B(w^p_{i+1/2,j})f_{i+1,j}
 \right)
 \right].
 \label{eq:radial_fv_flux}
\end{equation}
When $C_{A,i+1/2}=0$, Eq.~\eqref{eq:radial_fv_flux} is replaced by its first-order
upwind drift limit.

For an interior pitch face, define
\begin{equation}
 D_\xi(p,\xi)=\frac{\nu_D^{({\rm sel})}(p)}{2}(1-\xi^2),
\end{equation}
\begin{equation}
 A^\xi_{i,j+1/2}=A_\xi(p_i,\xi_{j+1/2}),\qquad
 D^\xi_{i,j+1/2}=D_\xi(p_i,\xi_{j+1/2}),
\end{equation}
and the face Peclet argument
\begin{equation}
 w^\xi_{i,j+1/2}=\frac{A^\xi_{i,j+1/2}d_{\xi,j+1/2}}{D^\xi_{i,j+1/2}}.
 \label{eq:pitch_peclet}
\end{equation}
For $D^\xi_{i,j+1/2}>0$, the integrated pitch face flux is
\begin{equation}
 F^\xi_{i,j+1/2}
 =\frac{2\pi}{3}
  \left(p_{i+1/2}^3-p_{i-1/2}^3\right)
 \left[
  \frac{D^\xi_{i,j+1/2}}{d_{\xi,j+1/2}}
 \left(
   \mathcal B(-w^\xi_{i,j+1/2})f_{ij}
  -\mathcal B(w^\xi_{i,j+1/2})f_{i,j+1}
 \right)
 \right].
 \label{eq:pitch_fv_flux}
\end{equation}
When $D^\xi_{i,j+1/2}=0$, Eq.~\eqref{eq:pitch_fv_flux} is replaced by its first-order
upwind drift limit.  The same Bernoulli/Chang--Cooper rule is therefore used in both
momentum coordinates.  Accuracy is obtained by mesh refinement rather than by
introducing a more elaborate spatial discretization at this stage.
The superscript ``sel'' emphasizes that the stencil is independent of the selected
small-angle collision model; only the precomputed coefficient values differ.

\subsection{Boundary conditions}

At the inner radial face, an origin-touching grid with $p_{\min}=0$ would use regularity
and the vanishing radial surface area to impose
\begin{equation}
 F^p_{1/2,j}=0.
 \label{eq:inner_zero_flux}
\end{equation}
For the current shifted domain $p_{\min}>0$, an absorbing zero-inflow/open-outflow
condition is used.  The exterior distribution is zero: advective inflow is suppressed,
advective drift with $A_p(p_{\min},\xi_j)<0$ leaves the represented domain, and, when
radial diffusion is enabled, the boundary value $f=0$ is imposed at the face using the
geometric center-to-face distance $d_{p,1/2}$.  This is the boundary used by the current
suprathermal configuration.
Numerically, the same Chang--Cooper face flux as the interior is used with the exterior
state set to zero; for $C_A>0$ this gives
\begin{equation}
 F^p_{1/2,j}=-2\pi p_{\min}^2\Delta\xi_j\,
 \frac{C_A}{d_{p,1/2}}\mathcal B(w^p_{1/2,j})f_{1j},
 \qquad
 w^p_{1/2,j}=\frac{A_p(p_{\min},\xi_j)d_{p,1/2}}{C_A}.
 \label{eq:inner_absorbing_flux}
\end{equation}

At $\xi=\pm1$, both the pitch drift and diffusion coefficients vanish through
$1-\xi^2$, so
\begin{equation}
 F^\xi_{i,1/2}=F^\xi_{i,N_\xi+1/2}=0.
 \label{eq:pitch_boundary_flux}
\end{equation}

The outer radial boundary is a zero-inflow, open-outflow boundary.  The first-order
boundary flux is
\begin{equation}
 F^p_{N_p+1/2,j}
 =2\pi p_{\max}^2\Delta\xi_j,
  \left[A_p(p_{\max},\xi_j)\right]_+ f_{N_p,j},
 \qquad [a]_+\equiv\max(a,0),
 \label{eq:outer_flux}
\end{equation}
with zero diffusive flux through the outer face and no incoming distribution supplied
from $p>p_{\max}$.  The value of $p_{\max}$ is chosen sufficiently large that the
reported solution is insensitive to further domain extension.  Domain-size convergence
is tested by increasing $p_{\max}$ while independently monitoring the geometric radial
resolution.

\subsection{Discrete Chiu--Harvey large-angle gain}

The large-angle configuration is selected independently of the small-angle collision
model.  It is either disabled, $L_{\rm LA}=0$, or the Chiu--Harvey gain is enabled.  For
constant prescribed plasma parameters the Chiu--Harvey source is a fixed linear map of
the evolving finite-volume state and is retained in the same fully implicit TR--BDF2
operator.

\paragraph{Chiu--Harvey finite-energy aligned-primary gain.}
The physical Chiu--Harvey gain operator is fully defined in
Sec.~1, in particular by Eqs.~\eqref{eq:ch_primary_spectrum},
\eqref{eq:ch_primary_root}, and \eqref{eq:ch_reduced_source}, together with the
knock-on kinematic and cutoff conditions stated there.  The numerical model does not
introduce a different large-angle collision model; it only evaluates that operator on
the finite-volume mesh.

For each radial cell, the pitch-integrated primary distribution appearing in
Eq.~\eqref{eq:ch_primary_spectrum} is evaluated directly from the current 2P
solution as
\begin{equation}
 \boxed{
 \mathcal F_i(\tau)
 =\Delta\xi\sum_{j=1}^{N_\xi} f_{ij}(\tau)
 }
 \label{eq:discrete_primary_spectrum}
\end{equation}
which is the pitch integral of the piecewise-constant reconstruction in radial cell
$i$.  Thus the primary distribution used by the Chiu--Harvey source is the
pitch-integrated distribution of the evolving kinetic solution itself; no separate
primary distribution or analytic closure is introduced.

For a target cell center $(p_i,\xi_j)$, evaluate the required primary momentum
$p'_0(p_i,\xi_j)$ from Eq.~\eqref{eq:ch_primary_root}.  If the root and the associated
secondary satisfy all of the knock-on kinematic and large-angle-cutoff conditions of
Sec.~1, and if $p'_0$ lies inside the represented radial domain, let $k$ be the radial
cell containing that momentum,
\begin{equation}
 p_{k-1/2}\le p'_0(p_i,\xi_j)<p_{k+1/2}.
 \label{eq:ch_primary_cell}
\end{equation}
The first-order piecewise-constant representation gives
\begin{equation}
 \mathcal F\!\left(p'_0,\tau\right)\approx\mathcal F_k(\tau).
 \label{eq:ch_piecewise_primary}
\end{equation}
The normalized Chiu--Harvey source at the target cell center is then obtained by direct
substitution of this value into Eq.~\eqref{eq:ch_reduced_source}, using the normalization
and event-rate scale of Sec.~\ref{sec:normalization}.  Denote that value by
\begin{equation}
 S_{{\rm CH},ij}(\tau)
 =S_{\rm CH}[f](p_i,\xi_j,\tau).
 \label{eq:ch_cell_center_source}
\end{equation}
If the Section-1 kinematic or cutoff conditions are not satisfied, then
$S_{{\rm CH},ij}=0$.

The corresponding finite-volume contribution is evaluated with the same
piecewise-constant, first-order cell representation used for the rest of the baseline
solver,
\begin{equation}
 \boxed{
 \left.\frac{dN_{ij}}{d\tau}\right|_{\rm CH}
 =V_{ij}S_{{\rm CH},ij}
 }
 \label{eq:discrete_ch_source}
\end{equation}
where $V_{ij}$ is given by Eq.~\eqref{eq:cellvolume}.  Therefore the discrete
large-angle calculation is simply
\begin{equation}
 f_{ij}
 \;\longrightarrow\;
 \mathcal F_i=\Delta\xi\sum_j f_{ij}
 \;\longrightarrow\;
 S_{{\rm CH},ij}
 \;\longrightarrow\;
 \left.dN_{ij}/d\tau\right|_{\rm CH}.
\end{equation}
No additional physical approximation is introduced in Sec.~2.  Spatial accuracy of this
first-order evaluation is established by refinement of the geometric $p$ and uniform
$\xi$ meshes.
For fixed physical parameters and a fixed mesh, Eqs.~\eqref{eq:discrete_primary_spectrum}--
\eqref{eq:discrete_ch_source} define a time-independent linear mapping of the evolving
cell contents.


\paragraph{Exact discrete pitch-reduction factorization.}
For the algebra used by the PETSc implementation it is convenient to introduce the radial
cell phase volume
\begin{equation}
 V_i^{(p)}=\frac{2\pi}{3}
 \left(p_{i+1/2}^3-p_{i-1/2}^3\right)
\end{equation}
and the radial particle content
\begin{equation}
 \boxed{
 M_i(\tau)\equiv\sum_{j=1}^{N_\xi}N_{ij}(\tau)
 =V_i^{(p)}\mathcal F_i(\tau).
 }
 \label{eq:ch_radial_mass}
\end{equation}
Let $R$ denote the fixed pitch-reduction map
\begin{equation}
 \bm M=R\bm N,
 \qquad
 (R\bm N)_i=\sum_{j=1}^{N_\xi}N_{ij}.
 \label{eq:ch_reduction_matrix}
\end{equation}
For every active Chiu--Harvey target cell $q\equiv(i,j)$, let $k(q)$ denote the
primary radial cell selected by Eq.~\eqref{eq:ch_primary_cell}.  All target geometry,
knock-on factors, normalization factors, target-cell volume, and the factor
$1/V_{k(q)}^{(p)}$ are time independent and may therefore be collected into a scalar
$c_q$.  The discrete source then has the exact form
\begin{equation}
 \left.\frac{dN_q}{d\tau}\right|_{\rm CH}
 =c_q M_{k(q)}.
 \label{eq:ch_mass_scatter}
\end{equation}
Define the fixed scatter map $U$ by $(U\bm M)_q=c_qM_{k(q)}$ on active targets and
zero otherwise.  The complete discrete Chiu--Harvey operator is therefore
\begin{equation}
 \boxed{L_{\rm CH}=UR.}
 \label{eq:ch_UR_factorization}
\end{equation}
Equations~\eqref{eq:ch_radial_mass}--\eqref{eq:ch_UR_factorization} are not a new
quadrature or a new physical approximation: they are an exact factorization of
Eqs.~\eqref{eq:discrete_primary_spectrum}--\eqref{eq:discrete_ch_source}.  In particular,
the source can be applied by one pitch reduction per radial shell followed by one scatter
to each active target instead of repeating the same pitch sum separately for every target.
The expanded and reduced source applications are required to agree to floating-point
accuracy.

\subsection{Semidiscrete system and discrete balance}

Flatten the one-based cell indices with
\begin{equation}
 k(i,j)=(i-1)N_\xi+j,
 \qquad k=1,\ldots,N_pN_\xi.
 \label{eq:flattening}
\end{equation}
After discretizing every local differential term, let $L_0$ denote the time-independent
finite-volume operator containing electric acceleration, the selected small-angle
collision coefficients, synchrotron reaction, and all boundary fluxes.  The selected
large-angle operator is
\begin{equation}
 L_{\rm LA}=\begin{cases}
 0, & \text{no large-angle gain},\\
 U_{\rm CH}R_{\rm CH}, & \text{Chiu--Harvey}.
 \end{cases}
 \label{eq:large_angle_discrete_choices}
\end{equation}
The complete semidiscrete equation is
\begin{equation}
 \boxed{
 L=L_0+L_{\rm LA},\qquad
 \frac{d\bm N}{d\tau}=L\bm N.
 }
 \label{eq:semidiscrete}
\end{equation}
For standalone mode this is one fixed linear operator; no operator splitting is
introduced.  In coupled mode the operator is rebuilt at each Newton iterate from
the current bulk stage state, while its finite-volume sparsity pattern remains fixed.

Summing the physical finite-volume equations cancels every shared interior conservative
flux exactly.  Writing $\Phi_{\rm in}\ge0$ and $\Phi_{\rm out}\ge0$ for particle loss
through the inner and outer radial boundaries, respectively, gives
\begin{equation}
 \boxed{
 \frac{d}{d\tau}\bm 1^T\bm N
 =-\Phi_{\rm in}-\Phi_{\rm out}+\bm 1^TL_{\rm LA}\bm N.
 }
 \label{eq:discrete_particle_balance}
\end{equation}
For the zero-flux production inner boundary $\Phi_{\rm in}=0$; it can be nonzero for the
shifted absorbing boundary.  For Chiu--Harvey the production term is
$\bm1^TU_{\rm CH}R_{\rm CH}\bm N$, and it vanishes when the gain is disabled.

\subsection{Implicit time stepping and planned TR--BDF2 method}

For fixed-step standalone mode, the normalized time grid is uniform,
\begin{equation}
 \tau^n=n\Delta\tau,\qquad n=0,1,\ldots,N_t.
 \label{eq:uniform_time_grid}
\end{equation}
Adaptive mode instead accepts a sequence of positive trial steps $h_n$ with
$\tau^{n+1}=\tau^n+h_n$.  The complete selected semidiscrete operator of
Eq.~\eqref{eq:semidiscrete} is advanced implicitly.  No enabled physical term is
lagged, split off, or treated explicitly.

The current constant-parameter CPU prototype uses fixed-step fully implicit TR--BDF2
with a shared factored stage matrix.  The electric-field coupled prototype still uses
variable-step fully implicit BDF2 with a backward-Euler comparison solve.  The TR--BDF2
method below is the target integrator for that coupled path as well.

The two-stage TR--BDF2 method uses
\begin{equation}
 d=1-\frac{1}{\sqrt{2}},\qquad h=\Delta\tau
 \quad\text{(or }h=h_n\text{ in adaptive mode)}.
 \label{eq:trbdf2_parameter}
\end{equation}
For a fixed operator $L$, define the common stage matrix
\begin{equation}
 \mathcal A_d(h)=I-dhL.
 \label{eq:trbdf2_stage_matrix}
\end{equation}
The first stage is the trapezoidal solution at the intermediate TR--BDF2 stage,
\begin{equation}
 \boxed{
 \mathcal A_d(h)\bm N^\gamma
 =\left(I+dhL\right)\bm N^n
 =2\bm N^n-\mathcal A_d(h)\bm N^n.
 }
 \label{eq:trbdf2_stage1}
\end{equation}
The endpoint stage uses the same matrix,
\begin{equation}
 \boxed{
 \mathcal A_d(h)\bm N^{n+1}
 =\bm N^n+\frac{h}{2\sqrt{2}}
   \left(L\bm N^n+L\bm N^\gamma\right).
 }
 \label{eq:trbdf2_stage2}
\end{equation}
Thus each accepted fixed-step advances through two sparse direct solves using the same
stage matrix.  The constant-parameter CPU path now implements these equations directly.

When adaptive stepping is enabled, the embedded error estimate solves the same stage
matrix with
\begin{equation}
 \mathcal A_d(h)\bm e^{n+1}
 =\bm N^n+h\left(
 b_0L\bm N^n+b_1L\bm N^\gamma+b_2L\bm N^{n+1}
 \right)-\bm N^{n+1},
 \label{eq:trbdf2_error}
\end{equation}
where
\begin{equation}
 (b_0,b_1,b_2)=
 (0.2154822031355755,\,0.6868867239266070,\,0.09763107293781750).
 \label{eq:trbdf2_error_weights}
\end{equation}
The scaled norm of $\bm e^{n+1}$ controls acceptance and the next trial step.  If a
trial step changes, the numerical matrix values are reassembled and the PETSc sparse
direct solver is refactorized while its CSR topology is reused.

With Chiu--Harvey enabled, introduce the radial auxiliary state
$\bm M=R_{\rm CH}\bm N$.  The stage matrix is the exact augmented system
\begin{equation}
 \boxed{
 \mathcal A_{d,\rm aug}(h)=
 \begin{bmatrix}
  I-dhL_0 & -dhU_{\rm CH}\\
  -R_{\rm CH} & I
 \end{bmatrix}.
 }
 \label{eq:trbdf2_augmented_matrix}
\end{equation}
The right-hand side is the corresponding physical TR--BDF2 stage right-hand side
with a zero auxiliary block.  Eliminating $\bm M$ recovers
$L=L_0+U_{\rm CH}R_{\rm CH}$ exactly.  The auxiliary constraint contains no time
derivative, so its identity block is unchanged between stages.
When Chiu--Harvey is disabled, the auxiliary state and constraint rows are omitted
entirely; the PETSc system then contains only the $N_pN_\xi$ physical cell contents.

In the planned coupled mode, each TR--BDF2 stage is one nonlinear system containing the kinetic
cell state, any Chiu--Harvey auxiliary state, the bulk temperature, every charge-state
population, the bulk current, and an explicit kinetic-current closure scalar.  The
electric field is algebraically eliminated through Ohm's law.  Let
\begin{equation}
 \bm x=(\bm N,\bm M,T_e,\{n_{a,q}\},I_p,I_{\rm RE})^T,
 \qquad
 \bm F_{\rm stage}(\bm x)=\bm 0
 \label{eq:coupled_stage_system}
\end{equation}
denote the assembled stage system, where $\bm M=R_{\rm CH}\bm N$ is present only
for the augmented Chiu--Harvey formulation.  Newton solves
\begin{equation}
 \bm J_{\rm stage}(\bm x_k)\,\delta\bm x_k
 =-\bm F_{\rm stage}(\bm x_k),\qquad
 \bm x_{k+1}=\bm x_k+\alpha_k\delta\bm x_k,
 \label{eq:coupled_newton_step}
\end{equation}
with analytic kinetic, kinetic--bulk, bulk--kinetic, and bulk Jacobian blocks.
The scalar kinetic-current closure is intentionally left open pending the current-model
decision described below.  Positivity damping is applied when required.  Newton first
converges the TR stage, then the BDF2
stage; only a converged pair is eligible for timestep error estimation and
acceptance.

\subsection{Sparse direct solution with PETSc}

The local physical block always has the same finite-volume stencil.  Changing between
the two small-angle collision models changes coefficient values but not the CSR topology.
The Chiu--Harvey configuration adds the exact low-dimensional radial auxiliary system of
Eq.~\eqref{eq:trbdf2_augmented_matrix}.

For one standalone simulation, the direct-solve procedure is:
\begin{enumerate}
 \item construct the geometric-radial finite-volume grid and cell geometry;
 \item project the prescribed initial distribution using Eq.~\eqref{eq:initial_projection};
 \item evaluate the selected small-angle collision coefficients once;
 \item if Chiu--Harvey is enabled, construct its fixed target geometry and the exact
       reduction/scatter maps $R_{\rm CH}$ and $U_{\rm CH}$;
 \item construct and validate the fixed CSR graph for the physical or augmented system;
 \item assemble the implicit TR--BDF2 stage matrix for fixed-parameter mode, using
       only the physical block when Chiu--Harvey is disabled;
 \item perform one sparse symbolic analysis of that fixed graph;
 \item numerically factor the stage matrix and solve the implicit stage;
 \item in adaptive mode, reassemble and refactor numerical values when the trial
       timestep changes while reusing the CSR topology;
 \item in coupled mode, construct the monolithic stage CSR graph once, assemble the
       analytic Newton Jacobian values at each iterate, and refactorize while reusing
       the topology.
\end{enumerate}
Fixed-parameter mode uses two TR--BDF2 stages with a shared matrix.  The current
electric-field adaptive prototype uses variable-step BDF2 with a backward-Euler
comparison solve for its local error estimate.  Coefficient or timestep changes require
numerical refactorization but not CSR symbolic reanalysis when topology remains unchanged.

For Chiu--Harvey, the exact augmented graph contains the local physical stencil, one
auxiliary-column coupling for each active target, and $N_p$ radial constraints
$M_i-\sum_jN_{ij}=0$.  Its raw graph size is $O(N_pN_\xi)$ and avoids the quadratic
pitch growth that would result from explicitly multiplying $U_{\rm CH}R_{\rm CH}$.

The CPU implementation uses PETSc distributed AIJ matrices and vectors with FP64
arithmetic.  CSR column indices, finite-volume slot maps, auxiliary couplings, and graph
validation are constructed during setup.  PETSc's configured sparse direct solver
performs numerical factorization and triangular solves.

For any solved linear system $A\bm x=\bm b$, the relative residual
\begin{equation}
 r_{\rm lin}=\frac{\|A\bm x-\bm b\|_2}{\|\bm b\|_2}
 \label{eq:linear_residual}
\end{equation}
is evaluated directly.  When the Chiu--Harvey augmented formulation is active, the
constraint defect $\|\bm M-R_{\rm CH}\bm N\|$ is checked independently.

\subsection{Discrete diagnostics and numerical qualification}

The solver reports direct moments of the represented kinetic distribution.  The
baseline current diagnostic is the total kinetic parallel current, not a
threshold-defined runaway current.  Its units are current density
$\mathrm{A\,m^{-2}}$; multiplying by the configured area gives total current
$\mathrm{A}$.

The normalized represented density is exactly
\begin{equation}
 \boxed{
 \frac{n_F}{n_{\rm ref}}=\sum_{i=1}^{N_p}\sum_{j=1}^{N_\xi}N_{ij}
 }
 \label{eq:discrete_density}
\end{equation}
because $N_{ij}$ is the particle content in the normalized phase-space measure.

The dimensional parallel current density corresponding to the piecewise-constant
reconstruction is
\begin{equation}
 \boxed{
 j_\parallel
 =-e n_{\rm ref}c
 \sum_{i=1}^{N_p}\sum_{j=1}^{N_\xi}
 f_{ij}\,W_i^{(v)}W_j^{(\xi)}
 }
 \label{eq:discrete_current}
\end{equation}
with exact cell weights
\begin{align}
 W_i^{(v)}
 &=2\pi\int_{p_{i-1/2}}^{p_{i+1/2}}
       p^2\frac{p}{\sqrt{1+p^2}}\,dp
 \nonumber\\
 &=2\pi\left[
       \frac{\gamma^3}{3}-\gamma
       \right]_{p_{i-1/2}}^{p_{i+1/2}},
 \qquad \gamma=\sqrt{1+p^2},
 \label{eq:current_radial_weight}\\
 W_j^{(\xi)}
 &=\int_{\xi_{j-1/2}}^{\xi_{j+1/2}}\xi\,d\xi
 =\frac{\xi_{j+1/2}^2-\xi_{j-1/2}^2}{2}.
 \label{eq:current_pitch_weight}
\end{align}
Thus neither density nor total kinetic current requires an additional diagnostic
quadrature.  A thresholded diagnostic uses the same weights with the indicator in
Eq.~\eqref{eq:threshold_runaway_current}, evaluated on the discrete cell energy.
The provisional CPU coupling uses that thresholded diagnostic in its bulk Ohm and
induction equations; the full-moment alternative remains under review.  In the current
prototype $P_m=P_{\min}$, so it is not an independent cutoff.

\paragraph{Numerical refinement parameters.}
Every selected physical configuration is qualified by independent refinement of
\begin{equation}
 \boxed{
 N_p,\qquad N_\xi,\qquad p_{\max},\qquad
 q_{\rm init},\qquad \Delta\tau
 }
 \label{eq:convergence_parameters}
\end{equation}
with $q_{\rm coll}$ added as an independent convergence parameter only for the
finite-temperature relativistic coefficient evaluation.  The fully-relativistic
asymptotically matched coefficients are analytic and use no collision quadrature.

Spatial refinement is performed on the current geometric-radial and uniform-pitch grid
family of Eqs.~\eqref{eq:geometric_p_grid}--\eqref{eq:current_xi_grid}.  Radial and
pitch resolution are tested independently by increasing $N_p$ and $N_\xi$.  Domain-size
convergence is a separate test: $p_{\max}$ is increased while the radial resolution is
also monitored, so that the physical domain and local momentum resolution are not
confounded.  The accepted $p_{\max}$ is
large enough that the represented distribution and reported moments are insensitive to
a further increase of the outer boundary over the simulated time interval.

For the finite-temperature model, $q_{\rm coll}$ is increased until the collision
coefficients are unchanged to the accuracy required by the mesh-converged kinetic
solution.  The fully-relativistic asymptotically matched model instead requires
asymptotic checks of its analytic direct evaluation.  In all configurations, $q_{\rm init}$ is increased
until the projected initial cell contents in Eq.~\eqref{eq:initial_projection} are
converged.  These checks are performed independently of mesh refinement.

\paragraph{Required numerical qualification tests.}
These tests establish that each implemented configuration discretizes its Section-1
operator correctly.  They are distinct from external literature benchmarks, which test
whether the resulting physical model reproduces established kinetic results.

\begin{enumerate}
 \item \textbf{Small-angle coefficient evaluation.}
       For the finite-temperature model, the direct scaled formulas and low-$p$ Taylor
       formulas must agree in an overlap interval and increasing $q_{\rm coll}$ must
       leave the coefficients unchanged.  For the fully-relativistic asymptotically matched model,
       Eqs.~\eqref{eq:matched_numeric_cf}--\eqref{eq:matched_numeric_nud} must reproduce the
       Section-1 operator in Eq.~\eqref{eq:matched_tp_operator} after coordinate
       transformation; the direct and small-$x$ evaluations of $\Psi$ must overlap, and
       the non-relativistic and relativistic asymptotic limits must be recovered.
       For both base models, the energy-dependent Coulomb-logarithm layer must satisfy
       Eqs.~\eqref{eq:common_log_numeric_cf}--\eqref{eq:common_log_numeric_nud}; setting
       $R_{ee}=R_{ei}=1$ must recover the thermal-log base exactly, and $C_A$ must remain
       unchanged.  Fully stripped ions must set only the partial-screening terms to zero,
       leaving the common Coulomb-logarithm corrections active.  With bound electrons
       present, the additional changes in $C_F$ and $\nu_D$ must equal
       Eqs.~\eqref{eq:screening_numeric_cf} and \eqref{eq:screening_numeric_nud},
       respectively, again with $C_A$ unchanged.

 \item \textbf{Finite-volume balance.}
       For arbitrary test states, summing the semidiscrete physical-cell equations must
       reproduce Eq.~\eqref{eq:discrete_particle_balance} to floating-point accuracy.
       Every interior face contribution must cancel pairwise between adjacent cells.

 \item \textbf{Collision equilibrium.}
       With electric field, synchrotron radiation, and large-angle gain removed, the
       appropriate Maxwell--J{\"u}ttner equilibrium at the prescribed background
       temperature must approach a stationary discrete solution as the mesh and
       any applicable coefficient quadrature are refined.  Exact preservation is not
       required from the first-order spatial scheme; the equilibrium defect must converge
       to zero.

 \item \textbf{Spatial convergence of the kinetic equation.}
       Smooth verification problems must recover first-order asymptotic convergence with
       respect to the first-order drift discretization as the geometric radial spacing
       and $\Delta\xi$ are independently reduced.  Production configurations must then
       be continued to finer meshes until
       the full distribution and reported moments are insensitive to further refinement.

 \item \textbf{Chiu--Harvey discretization, when selected.}
       For a prescribed smooth two-dimensional distribution, the discrete primary pitch
       integral and the resulting cell-centered source must converge to the Section-1
       Chiu--Harvey expressions under independent radial and $\Delta\xi$ refinement.
       The test includes the finite-energy primary root, kinematic support, and fixed
       large-angle cutoff.

 \item \textbf{Chiu--Harvey algebraic factorization, when selected.}
       On a fixed mesh and state, the pitch-reduction/scatter application must agree with
       the explicitly expanded $N_\xi$-term source sum to floating-point accuracy.  The
       augmented TR--BDF2 stages must reproduce the physical solution of the expanded
       stages to the accuracy permitted by the sparse residual, and
       $\bm M=R_{\rm CH}\bm N$ must hold to floating-point accuracy.

 \item \textbf{Outer-domain convergence.}
       At fixed local radial resolution and $\Delta\xi$, increasing $p_{\max}$ must leave the solution
       on the original momentum interval and the reported moments unchanged to the
       required accuracy.

 \item \textbf{Temporal convergence.}
       With a spatially converged mesh, reducing $\Delta\tau$ must recover second-order
       convergence of the fully implicit TR--BDF2 evolution over a fixed physical
       interval.  Adaptive runs must also demonstrate error-controlled step acceptance
       and stable convergence under timestep refinement.

 \item \textbf{Sparse direct solve.}
       Every TR--BDF2 stage solve and adaptive error-estimate solve must satisfy
       Eq.~\eqref{eq:linear_residual}; the residual must remain small compared with the
       demonstrated spatial and temporal discretization errors.  Any active auxiliary
       constraint must satisfy its separate defect check.
\end{enumerate}

Unresolved coarse meshes may be used for algebraic, indexing, and implementation tests,
but they are not evidence of physical convergence.  The numerical strategy is deliberately to use a simple first-order finite-volume
discretization on a sufficiently resolved CPU mesh and to demonstrate convergence
directly for each selected physical configuration.


%===============================================================================
\section{Bulk-plasma coupling: current prototype and target}
\label{sec:bulk_coupling}

Standalone mode advances the kinetic distribution using a prescribed, constant plasma
state.  Coupled operation uses the analytic-Newton prototype with PETSc linear solves
and host-side bulk residual evaluation.

\subsection{Bulk state and collisional-radiative equations}

For each coupled atomic species $a$, let $Z_a$ be its atomic number and let
$n_{a,q}$ denote the density in charge state $q$, with $q=0,\ldots,Z_a$.  The total
nuclei density evolves through the neutral injection source,
\begin{equation}
 \sum_{q=0}^{Z_a} n_{a,q}=n_{a,\rm tot}(t),
 \qquad
 n_{a,\rm tot}(t)=n_{a,\rm base}
 +\Delta n_a\int_0^t
 \frac{\mathbf 1_{(t_a,t_a+\Delta t_a)}(s)}{\Delta t_a}\,ds.
\label{eq:bulk_nuclei_conservation}
\end{equation}
The implemented initial state places each species in its configured initial charge state
$q_{a,0}$,
\begin{equation}
 n_{a,q}(0)=n_{a,\rm tot}(0)\,\delta_{q,q_{a,0}}.
 \label{eq:bulk_initial_charge_state}
\end{equation}
The free-electron density and effective charge are computed from the charge-state
populations,
\begin{align}
 n_e &= \sum_a\sum_{q=0}^{Z_a}q\,n_{a,q},
 & Z_{\rm eff} &=
 \frac{\displaystyle\sum_a\sum_{q=0}^{Z_a}q^2 n_{a,q}}{n_e}.
 \label{eq:bulk_electron_density_zeff}
\end{align}
Thus a background deuterium plasma with $n_e=10^{20}\,\mathrm{m}^{-3}$ is represented
by a hydrogen OpenADAS bundle mapped to deuterium and a fully stripped $q=1$ population;
additional D, Ne, or Ar populations add their own free electrons through the same equation.
The cold target density used by the large-angle source is
\begin{equation}
 n_t=n_e+\sum_a\sum_{q=0}^{Z_a}(Z_a-q)n_{a,q},
 \label{eq:bulk_target_density}
\end{equation}
unless standalone input supplies an explicit $n_t$ override.

OpenADAS ADF11 SCD and ACD coefficients are evaluated at the current $(T_e,n_e)$ in native
$\mathrm{cm^3\,s^{-1}}$ units.  With the ADF11 charge-block convention used by the code,
define
\begin{align}
 S_{a,q}(T_e,n_e) &= \mathrm{SCD}_{a,q+1}(T_e,n_e),
 &&q=0,\ldots,Z_a-1,\\
 \alpha_{a,q}(T_e,n_e) &= \mathrm{ACD}_{a,q}(T_e,n_e),
 &&q=1,\ldots,Z_a.
 \label{eq:bulk_adas_rates}
\end{align}
The ECD table supplies a charge-state mean ionisation potential
$I^{\rm ECD}_{a,q}(T_e,n_e)$ in eV.  The reader interpolates ECD values directly in
the tabulated values, rather than in logarithmic form.  When explicit state-resolved
excitation energies are absent, the coupled kinetic driver uses ECD values for the
$I_{a,q}$ inputs in Eq.~\eqref{eq:physical_h}.  This requires verifying that the ECD
definition matches the Hesslow mean-excitation definition for the selected data set.
The ADF11 block $Z1=q+1$ supplies the value for charge state $q$; the bare state has
no next ionisation transition and uses zero.
ECD values are not successive ionization energies $I_{a,r}$ in
Eq.~\eqref{eq:bulk_energies}; the current bulk energy ledger keeps those transition
energies as explicit species inputs.
The collisional-radiative population equations are
\begin{equation}
 \frac{d n_{a,q}}{dt}
 = n_e\left[S_{a,q-1}n_{a,q-1}+\alpha_{a,q+1}n_{a,q+1}
 -(S_{a,q}+\alpha_{a,q})n_{a,q}\right]
 +\mathcal I_{a,q},
 \label{eq:bulk_charge_state_rhs}
\end{equation}
where terms with invalid charge indices are zero.  In the nonlinear stage system, one
redundant bare-state equation per species is replaced by the exact constraint
Eq.~\eqref{eq:bulk_nuclei_conservation}; no post-root nuclei projection is applied.
The factors $n_e$ are converted to
$\mathrm{cm^{-3}}$ before multiplying the native OpenADAS coefficients.  With no injection,
the collisional-radiative terms conserve Eq.~\eqref{eq:bulk_nuclei_conservation}; with
injection, the source in Eq.~\eqref{eq:bulk_neutral_injection_source} sets the changing
species total.

The OpenADAS PLT and PRB coefficients provide line-plus-continuum and recombination-plus-
bremsstrahlung power.  In native $\mathrm{W\,cm^3}$ units, the implemented radiated power
density is
\begin{equation}
 P_{\rm rad}=10^{-6}n_e\sum_a\sum_{q=0}^{Z_a-1}
 n_{a,q}\left[\mathrm{PLT}_{a,q+1}(T_e,n_e)+\mathrm{PRB}_{a,q+1}(T_e,n_e)\right].
 \label{eq:bulk_radiated_power}
\end{equation}
The ADF11 tables use bilinear interpolation in
$x=\log_{10}(T_e/\mathrm{eV})$ and $y=\log_{10}(n_e/\mathrm{cm^{-3}})$, with the
coefficient interpolated in logarithmic form:
\begin{align}
 \log_{10}K(x,y)&=(1-f_x)(1-f_y)\log_{10}K_{00}
 +(1-f_x)f_y\log_{10}K_{01}\\
 &\quad+f_x(1-f_y)\log_{10}K_{10}+f_xf_y\log_{10}K_{11},\\
 K(x,y)&=10^{\log_{10}K(x,y)}.
 \label{eq:bulk_adas_interpolation}
\end{align}
Here $f_x$ and $f_y$ are fractional positions inside the enclosing native grid cell.
Interpolation occurs only inside the native temperature and density rectangle.  Missing
tables or out-of-range values are run-time errors.
The analytic Newton Jacobian differentiates this interpolation rule with respect to
$T_e$ and $n_e$.  ECD uses the same bilinear grid coordinates but interpolates the
tabulated mean-ionisation-potential values directly, with derivatives of that direct
interpolation rule.

\subsection{Bulk energy, resistivity, and current equations}

The implemented bulk energy state contains electron thermal energy and ionization energy,
\begin{align}
 U_e &= \frac32 e n_e T_e,\\
 U_{\rm at} &= e\sum_a\sum_{q=0}^{Z_a}E_{a,q}n_{a,q},
 & E_{a,q}&=\sum_{r=1}^{q}I_{a,r},\qquad E_{a,0}=0.
 \label{eq:bulk_energies}
\end{align}
Here $T_e$ is stored in electron-volts, $I_{a,r}$ is the ionization energy for transition
$r-1\to r$, and $eT_e$ converts electron-volts to joules.  The temperature equation is
\begin{equation}
 \frac{dT_e}{dt}
 =\frac{P_\Omega-P_{\rm rad}-dU_{\rm at}/dt
 -(3/2)eT_e\,dn_e/dt}{(3/2)e n_e}.
 \label{eq:bulk_temperature_rhs}
\end{equation}
The population equations determine the two derivatives in this expression,
\begin{align}
 \frac{dn_e}{dt}&=\sum_a\sum_{q=0}^{Z_a}q\frac{dn_{a,q}}{dt},\\
 \frac{dU_{\rm at}}{dt}&=e\sum_a\sum_{q=0}^{Z_a}E_{a,q}\frac{dn_{a,q}}{dt}.
 \label{eq:bulk_energy_derivatives}
\end{align}
The code includes no separate ion-temperature or transport-energy equation.

The thermal Coulomb logarithm and Spitzer-like resistivity used by the bulk model are
\begin{align}
 \ln\Lambda_0&=14.9-\frac12\ln\left(\frac{n_e}{10^{20}\,\mathrm{m}^{-3}}\right)
 +\ln\left(\frac{T_e}{10^3\,\mathrm{eV}}\right),\\
 \eta&=1.65\times10^{-9}\,\frac{Z_{\rm eff}\ln\Lambda_0}
 {(T_e/10^3\,\mathrm{eV})^{3/2}}\;\Omega\,\mathrm{m}.
 \label{eq:bulk_resistivity}
\end{align}
Let $I_p$ be the bulk current, $A$ the configured cross-sectional area, $R_0$ the
major radius, $L$ the lumped inductance, and $V_{\rm app}$ the applied voltage.  For
the provisional CPU closure, $j_{\rm RE}^{(\kappa)}$ is the thresholded non-thermal
kinetic current density from Eq.~\eqref{eq:threshold_runaway_current}.  The Ohm and
induction equations are
\begin{align}
 j_p&=\frac{I_p}{A},
 &I_{\rm RE}&=A j_{\rm RE}^{(\kappa)},\\
 E_\parallel&=\eta\left(j_p-j_{\rm RE}^{(\kappa)}\right),
 &P_\Omega&=E_\parallel j_p,\\
 \frac{dI_p}{dt}&=\frac{V_{\rm app}-R_p(I_p-I_{\rm RE})}{L},
 &R_p&=\frac{2\pi R_0\eta}{A}.
 \label{eq:bulk_ohm_induction}
\end{align}
The factor $A$ is used only to convert the current density to the lumped current
needed by the induction equation; $j_{\rm RE}^{(\kappa)}$ is the quantity used in
Ohm's law.  Whether this thresholded closure or the full $j_\parallel$ is the final
physical coupling choice remains open for a later qualification study.  The configured geometry is
\begin{equation}
 R_0=6\,\mathrm{m},\qquad a=2\,\mathrm{m},\qquad
 A=\pi a^2,\qquad L=\mu_0R_0.
 \label{eq:bulk_geometry}
\end{equation}
The magnetic field $B_T$ remains prescribed and is not exchanged with the bulk state.

The implemented homogeneous material-injection source is neutral and finite-duration.
For species $a$, with injected amount $\Delta n_a$, start time $t_a$, and duration
$\Delta t_a$, it is
\begin{equation}
 \mathcal I_{a,q}(t)=
 \frac{\Delta n_a}{\Delta t_a}
 \mathbf 1_{(t_a,t_a+\Delta t_a)}(t)\,\delta_{q0}.
 \label{eq:bulk_neutral_injection_source}
\end{equation}
The source is spatially homogeneous because the model has no spatial coordinate.  Its
particle, free-electron, and atomic-energy balances are
\begin{equation}
 \sum_q\mathcal I_{a,q}=\frac{dn_{a,\rm tot}}{dt},
 \qquad
 \left.\frac{dn_e}{dt}\right|_{\rm inj}=\sum_{a,q}q\mathcal I_{a,q},
 \qquad
 \left.\frac{dU_{\rm at}}{dt}\right|_{\rm inj}
 =e\sum_{a,q}E_{a,q}\mathcal I_{a,q}.
 \label{eq:bulk_injection_sources}
\end{equation}
The single-temperature model assigns no separate kinetic-energy reservoir to injected
neutrals.  Consequently $\mathcal I_{a,0}$ contributes no direct injection-energy term;
electron energy changes only through the existing ionization-energy and radiation terms.

\subsection{Bulk TR--BDF2 stages and kinetic exchange}

Write the bulk state as
\begin{equation}
 \bm y=(T_e,n_{1,0},\ldots,n_{a,q},\ldots,I_p)^T,
 \qquad \frac{d\bm y}{dt}=\bm f(\bm y,t;I_{\rm RE}).
 \label{eq:bulk_state_rhs}
\end{equation}
For each attempted step of size $h$, the host model solves the same TR--BDF2 stages as
the kinetic model, with $d=1-1/\sqrt2$ and $w=\sqrt2/4$,
\begin{align}
 \bm y^\gamma-\bm y^n
 &=dh\left[\bm f(\bm y^n,t_n;I_{\rm RE}^n)
 +\bm f(\bm y^\gamma,t_n+dh;I_{\rm RE}^\gamma)\right],\\
 \bm y^{n+1}-\bm y^n
 &=dh\,\bm f(\bm y^{n+1},t_n+h;I_{\rm RE}^{n+1})
 +hw\left[\bm f(\bm y^n,t_n;I_{\rm RE}^n)
 +\bm f(\bm y^\gamma,t_n+dh;I_{\rm RE}^\gamma)\right].
 \label{eq:bulk_trbdf2}
\end{align}
The monolithic driver solves the coupled kinetic and bulk stage equations together.
The provisional kinetic-current closure supplies $I_{\rm RE}=A j_{\rm RE}^{(\kappa)}$
from the non-thermal discrete moment.  Newton will converge the TR stage first and the
BDF2 stage second once TR--BDF2 is implemented; only then should the driver evaluate
the adaptive error estimator and accept the timestep.

At each kinetic stage, the bulk state maps to kinetic coefficients through
\begin{align}
 T_e^{\rm kin}&=T_e, & n_e^{\rm kin}&=n_e,
 &E_\parallel^{\rm kin}&=\eta\left(\frac{I_p-I_{\rm RE}}{A}\right),\\
 Z_{0a}^{\rm kin}&=0, & n_{a,q}^{\rm kin}&=n_{a,q},
 & (I_{a,q}^{\rm kin},\bar a_{a,q}^{\rm kin})&=(I_{a,q},\bar a_{a,q}).
 \label{eq:bulk_to_kinetic_mapping}
\end{align}
For a fully stripped species with no bound-electron correction, the driver instead passes
$Z_{0a}^{\rm kin}=Z_a$ and its current total density.  The kinetic driver rebuilds
state-resolved collision coefficients from these values and retains the prescribed $B_T$.
The production stage convergence test uses separate scaled nonlinear residual and
Newton-update norms.  It does not use the PETSc linear residual as a nonlinear
stopping criterion.  The bulk scaling uses $s_T$, $s_n$, and $s_I$ for temperature,
charge-state populations, and current; the kinetic rows use a configured kinetic
residual scale.  Each species also satisfies Eq.~\eqref{eq:bulk_nuclei_conservation}
inside the Newton residual.  Rejected adaptive trials restore the last accepted
kinetic and bulk states exactly before reducing the trial step.

The legacy fixed-point residual used by the Picard comparison driver is retained only
for regression studies and is not the production coupling equation.

The implemented coupling contract is:
\begin{itemize}
  \item the host bulk state contains electron temperature, charge-state populations for
        every coupled species, and total current; electron density and parallel electric
        field are derived by the bulk model.  The magnetic field remains the prescribed
        $B_T$ configuration parameter and is not currently exchanged;
  \item each configured injected species receives the homogeneous neutral source in
        Eq.~\eqref{eq:bulk_neutral_injection_source}; the source changes species nuclei
        density but adds no separate injection-energy state;
  \item the planned coupled timestep solves its TR and BDF2 stages with the monolithic
        analytic-Newton system in Eq.~\eqref{eq:coupled_stage_system}; each Newton
        iteration updates numerical CSR values and uses PETSc factorization or
        refactorization while retaining the fixed topology;
  \item the bulk equations include full charge-state populations, neutral finite-duration
        injection, nuclei-total constraints, state-resolved quasineutrality, energy,
        Ohm, and lumped-induction terms;
  \item particle/current, bulk-energy, collisional-radiative, Ohmic, and lumped-induction
        terms are evaluated by the implemented host bulk model and reported through
        stage diagnostics; and
  \item local OpenADAS ADF11 bundles and authoritative state-resolved excitation and
        screening data are required for each coupled species.  The maintained Newton
        report records the input TOML, stage diagnostics, bulk state, charge populations,
        current, timing, and memory; the standalone output path additionally provides
        restart/provenance payloads.
\end{itemize}

The analytic-Newton path is the intended coupled implementation.  Qualification records
must include nuclei-total conservation, analytic directional derivatives, Newton/Picard
equation equivalence, adaptive rejection/retry, topology reuse, and restart equivalence.
The final choice between the thresholded $j_{\rm RE}^{(\kappa)}$ closure and the full
$j_\parallel$ closure remains an explicit physics decision; it must not be hidden in a
diagnostic or inferred from the lumped current.

\begin{thebibliography}{99}

\bibitem{ConnorHastie1975}
J.~W. Connor and R.~J. Hastie,
``Relativistic limitations on runaway electrons,''
\emph{Nucl. Fusion} \textbf{15}, 415 (1975).

\bibitem{BraamsKarney1989}
B.~J. Braams and C.~F.~F. Karney,
``Conductivity of a relativistic plasma,''
\emph{Phys. Fluids B} \textbf{1}, 1355 (1989).

\bibitem{PikeRose2014}
O.~J. Pike and S.~J. Rose,
``Dynamical friction in a relativistic plasma,''
\emph{Phys. Rev. E} \textbf{89}, 053107 (2014).

\bibitem{Hesslow2017Screening}
L.~Hesslow, O.~Embr\'eus, A.~Stahl, T.~C. DuBois, G.~Papp,
S.~L. Newton, and T.~F{\"u}l{\"o}p,
``Effect of partially screened nuclei on fast-electron dynamics,''
\emph{Phys. Rev. Lett.} \textbf{118}, 255001 (2017).

\bibitem{Hesslow2018Collision}
L.~Hesslow, O.~Embr\'eus, M.~Hoppe, T.~C. DuBois, G.~Papp,
M.~Rahm, and T.~F{\"u}l{\"o}p,
``Generalized collision operator for fast electrons interacting with partially ionized impurities,''
\emph{J. Plasma Phys.} \textbf{84}, 905840605 (2018).

\bibitem{Andersson2001}
F.~Andersson, P.~Helander, and L.-G.~Eriksson,
``Damping of relativistic electron beams by synchrotron radiation,''
\emph{Phys. Plasmas} \textbf{8}, 5221 (2001).

\bibitem{Stahl2016}
A.~Stahl, O.~Embr{\'e}us, G.~Papp, M.~Landreman, and T.~F{\"u}l{\"o}p,
``Kinetic modelling of runaway electrons in dynamic scenarios,''
\emph{Nucl. Fusion} \textbf{56}, 112009 (2016).

\bibitem{Moller1932}
C.~M{\o}ller,
``Zur Theorie des Durchgangs schneller Elektronen durch Materie,''
\emph{Ann. Phys.} \textbf{406}, 531 (1932).

\bibitem{Chiu1998}
S.~C. Chiu, M.~N. Rosenbluth, R.~W. Harvey, and V.~S. Chan,
\emph{Nucl. Fusion} \textbf{38}, 1711 (1998).

\bibitem{Harvey2000}
R.~W. Harvey, V.~S. Chan, S.~C. Chiu, T.~E. Evans,
M.~N. Rosenbluth, and D.~G. Whyte,
\emph{Phys. Plasmas} \textbf{7}, 4590 (2000).


\end{thebibliography}

\end{document}
